\documentclass[11pt,a4paper]{article}
\usepackage[margin=25mm,headheight=15pt]{geometry}
\usepackage[T1]{fontenc}
\usepackage{amsmath,amsthm,mathtools}
\usepackage{newtxtext,newtxmath}
\usepackage{microtype,booktabs,array,longtable,enumitem}
\usepackage{xcolor,fancyhdr,needspace}
\usepackage[most]{tcolorbox}
\usepackage{xurl}
\usepackage[colorlinks=true,linkcolor=blue!42!black,citecolor=blue!42!black,urlcolor=blue!42!black]{hyperref}
\usepackage[nameinlink,noabbrev]{cleveref}
\numberwithin{equation}{section}
\newtheorem{theorem}{Theorem}[section]
\newtheorem{proposition}[theorem]{Proposition}
\newtheorem{lemma}[theorem]{Lemma}
\newtheorem{corollary}[theorem]{Corollary}
\theoremstyle{definition}
\newtheorem{definition}[theorem]{Definition}
\newtheorem{problem}{Research question}
\theoremstyle{remark}
\newtheorem{remark}[theorem]{Remark}
\newcommand{\C}{\mathbb C}
\newcommand{\R}{\mathbb R}
\newcommand{\N}{\mathbb N}
\newcommand{\Z}{\mathbb Z}
\newcommand{\ii}{\mathrm i}
\newcommand{\e}{\mathrm e}
\newcommand{\dd}{\,\mathrm d}
\newcommand{\B}{\mathcal B}
\newcommand{\cL}{\mathcal L}
\newcommand{\cE}{\mathcal E}
\newcommand{\PV}{\operatorname{PV}}
\newcommand{\Res}{\operatorname*{Res}}
\newcommand{\id}{\operatorname{id}}
\newcommand{\eps}{\varepsilon}
\newcommand{\wt}{\widetilde}
\newcommand{\norm}[1]{\left\lVert#1\right\rVert}
\newcommand{\abs}[1]{\left|#1\right|}
\DeclareMathOperator{\Ei}{Ei}
\setlist{itemsep=3pt,topsep=5pt}
\setlength{\parskip}{3.5pt}
\setlength{\parindent}{15pt}
\setlength{\emergencystretch}{2em}
\pagestyle{fancy}
\fancyhf{}
\fancyhead[L]{\small Exact $q$-product completion}
\fancyhead[R]{\small Stokes-aware reversion}
\fancyfoot[C]{\thepage}
\hypersetup{pdftitle={Exact q-Product Completion and Stokes-Aware Transseries Reversion},pdfauthor={Research report prepared for Vladimir Reshetnikov},pdfsubject={Mellin-Borel completion, reflection kernels, nonlinear inversion, q-gamma critical points}}

\begin{document}
\hypersetup{pageanchor=false}
\begin{titlepage}
\thispagestyle{empty}
\vspace*{10mm}
{\sffamily\small RESEARCH ARTICLE\hfill 4 OCTOBER 2026\par}
\vspace{19mm}
{\Huge\bfseries Exact $q$-Product Completion\par}
\vspace{4mm}
{\LARGE\bfseries and Stokes-Aware\par Transseries Reversion\par}
\vspace{10mm}
{\Large Reflection kernels, inverse exponential actions,\par
and the critical fold of the $q$-gamma function\par}
\vspace{15mm}
{\large Prepared for Vladimir Reshetnikov\par}
{\normalsize A continuation of the ProveIt transseries program\par}
\vspace{14mm}
\begin{tcolorbox}[colback=blue!3,colframe=blue!35!black,boxrule=.5pt,arc=1mm]
\textbf{Main result.} A Mellin--Borel argument gives the exact exponentially small completion of a shifted infinite $q$-product, including its jump-free part. A reflection criterion characterizes precisely when the completion vanishes. Transporting this data through regular and critical inverse charts yields convergent exponential-sector expansions and a uniform, action-halving $q$-gamma fold.
\end{tcolorbox}
\vfill
\noindent\textbf{Proof and novelty status.} This article contains conventional proofs and reproducible symbolic and high-precision checks, not a Lean formalization or an independent referee report. The scalar completion identity specializes to Conjecture~5 in Fantini--Rella, arXiv:2506.08265v2. General resurgent closure, Lagrange inversion, and the analytic quadratic normal form are credited as established tools. A first-publication priority claim is not made.
\end{titlepage}
\hypersetup{pageanchor=true}

\begin{abstract}
Complex transseries reversion must transport more than the formal power--logarithmic expansion: it must also transport the selected exponentially small completion, the Stokes transition maps, and the inverse sheet. For
\[
 G_a(t)=\log(\e^{-at};\e^{-t})_\infty,\qquad 0<a<1,\quad \Re t>0,
\]
we prove the exact formula
\[
 G_a=P_a+L_0R_a+\tfrac12\log(\e^{2\pi\ii a}Q;Q)_\infty
                    +\tfrac12\log(\e^{-2\pi\ii a}Q;Q)_\infty,
 \qquad Q=\e^{-4\pi^2/t},
\]
where $P_a$ is elementary and $L_0$ is the symmetric lateral Borel sum of an explicit Bernoulli series. A Mellin contour shift and a cotangent-kernel identity prove the equality, not merely its asymptotic expansion. The corresponding lateral formulas give a direct proof of a scalar summability identity formulated as a conjecture in the recent literature.

For arbitrary finite shift weights, exact reconstruction by the linear median is equivalent to antisymmetry under $a\mapsto1-a$, whereas absence of a Stokes jump is equivalent to symmetry. A finite-phase bound limits how many initial exponential actions can cancel. We then develop chart-level formal and analytic reversion, with explicit all-orders completion transport and an exact inverse Stokes cocycle. At a regular $q$-product cusp, the fundamental action becomes $24$ in the logarithmic target coordinate, with computable power--logarithmic dressing. At the minimum of $\log\Gamma_{\e^{-t}}(a)$, we locate the exponentially shifted critical point and value and prove a uniform two-scale inverse expansion. The forward action $4\pi^2/t$ becomes $2\pi^2/t$ in the critical layer. Convergent quadratic normal forms control the branch point itself, where ordinary inverse perturbation estimates are nonuniform. The results are accompanied by verification code, a source-status audit, and a research agenda.
\end{abstract}
\begingroup
\small
\setlength{\parskip}{0pt}
\tableofcontents
\endgroup
\newpage

\section{Research objective, contribution, and provenance}
\label{sec:introduction}

\subsection{What is being reversed?}
Throughout, \emph{reversion} means compositional inversion, not taking the multiplicative reciprocal of a series. There are three different inversion tasks. A formal equation may have a unique solution in a specified complete transseries algebra. A selected analytic realization may have an inverse on a sectorial chart. Finally, analytic continuation may move between inverse sheets. These tasks interact but are not equivalent.

The central issue in this article is a concrete one. Suppose an exact $q$-function is the sum of a Borel-resummed carrier and an exponentially small completion. What does the completion do to the inverse, particularly when the carrier derivative is small? A Stokes jump alone does not determine the completion: a nonzero flat contribution can have no jump at all. Near a critical value, omitting that contribution can change the leading scale of the inverse.

Our principal example is the shifted infinite product $G_a$. The shift variable $a$ is initially real and lies in $(0,1)$; the asymptotic variable $t$ is complex. A holomorphic parameter extension and an exact shift relation subsequently allow inversion in $a$ near the positive minimum of the $q$-gamma function.

\subsection{Precisely stated contributions}
The technical center is \cref{thm:completion}, proved by an explicit Mellin transform rather than inferred from residues. It includes a full description of the jump-free completion. \Cref{thm:reflection} turns that description into a necessary-and-sufficient reflection criterion for arbitrary finite weights. \Cref{prop:delay} bounds the first surviving action using only the number of distinct phases.

The inversion framework in \cref{sec:inverse} is deliberately chart-level. Its formal fixed-point lemma and Lagrange--B\"urmann formula are established mechanisms, included with proofs to make their hypotheses and domains visible. The application in \cref{sec:cusp} supplies explicit one- and two-exponential inverse coefficients and identifies the logarithmic target action. The critical-point results in \cref{sec:fold} supply the $q$-specific data that a generic fold theorem does not provide: the exact completion, the displaced discriminant, and a uniform inverse expansion in the layer where regular inversion ceases to be valid.

\subsection{Relation to the repository}
The repository comparison is pinned to ProveIt commit
\begin{center}
\texttt{b484725c893a7a9bde856151988996aef6965cb0}.
\end{center}
The consolidated transseries source is in
\begin{quote}\small
\path{Analysis/Transseries/docs/series-and-transseries/Transseries_And_Inversion/}.
\end{quote}
The separate article \emph{Nonlinear Stokes Transport under Logarithmic Inversion} already proves an all-orders inverse-transport formula, treats exact logarithmic cores, distinguishes arithmetic averaging from nonlinear reconstruction, and analyzes a Lambert-type fold \cite{repo,repotransport}. The Gaussian-binomial crossover work also distinguishes exact exponential completion from a truncation-error threshold \cite{repocrossover}.

Accordingly, neither the existence of nonlinear inverse sectors nor the generic square-root mechanism is claimed here as a new discovery. The extension is an exact shifted-product reconstruction theorem, a finite-weight classification, and their application to the actual $q$-gamma critical discriminant. The repository comparison is a focused statement-level comparison, not a line-by-line audit of its entire large corpus.

\subsection{Literature and the boundary of a conjecture claim}
Sauzin proves closure results for resurgent functions under nonlinear substitution, composition, functional inversion, and implicit equations \cite{sauzin}. Garoufalidis--Kashaev explicitly analyze Borel singularities and Stokes phenomena for Faddeev's quantum dilogarithm \cite{gk}. These are background, not claims of this article.

Fantini--Rella's version dated 1 April 2026 computes the relevant $q$-product Borel transforms and states a scalar exact summability formula as Conjecture~5 \cite{fr}. \Cref{cor:fr} gives a direct proof of that identity, with all logarithmic constants and contour orientations specified. The finite-weight criterion gives its odd-character median-reconstruction consequence. More recent work includes vector-valued modular resurgence and character-twisted Lambert series \cite{vectors,bd}; we do not assert that the much broader conjectures or phase prescriptions in those settings are settled here.

The source-status statement is version-specific. The proof below can be assessed independently of whether an equivalent modular identity has appeared in another normalization. An exhaustive priority assessment has not been completed. In particular, the standard eta transformation and Fourier-functional-equation machinery are not newly discovered identities.

\section{Conventions for sectors, summation, and transseries}
\label{sec:conventions}

\subsection{Domains and logarithms}
Put
\begin{equation}
 A=4\pi^2,\qquad A_0=\frac{\pi^2}{6},\qquad
 \mathbb H_R=\{t\in\C:\Re t>0\},\qquad Q(t)=\e^{-A/t}.
 \label{eq:constants}
\end{equation}
On $\mathbb H_R$ we take $-\pi/2<\arg t<\pi/2$. The function $Q$ satisfies $|Q|<1$. Uniform asymptotic assertions will be made on proper closed subsectors
\begin{equation}
 S_{\eps,T}=\{0<|t|<T:\ |\arg t|\leq\pi/2-\eps\},
 \qquad 0<\eps<\pi/2.
 \label{eq:sector}
\end{equation}
There $Q$ is uniformly exponentially small as $t\to0$.

Logarithms of infinite products are defined by their convergent logarithmic series. In particular,
\begin{equation}
 G_a(t)=-\sum_{m\geq1}\frac{\e^{-amt}}{m(1-\e^{-mt})}
 \label{eq:G}
\end{equation}
for real $a>0$ and $t\in\mathbb H_R$. This convention is not the instruction to take a principal logarithm of the already multiplied product; the two can differ by a multiple of $2\pi\ii$ after continuation. For complex $a$ we use a local continuation on a product neighborhood where the logarithmic series converges, and then the stated shift identities.

\subsection{Borel and lateral conventions}
For a series without constant term,
\begin{equation}
 \B\left[\sum_{n\geq1}c_nt^n\right](\xi)
 =\sum_{n\geq1}\frac{c_n\xi^{n-1}}{(n-1)!}.
 \label{eq:borel}
\end{equation}
Our Laplace integral has no extra factor $1/t$. The contours $L_+$ and $L_-$ go from zero to positive infinity, passing respectively above and below each positive real Borel pole. Upper small semicircles are traversed clockwise. Thus a pole with residue $r$ at $\omega>0$ contributes
\begin{equation}
 L_+-L_-=-2\pi\ii\,r\e^{-\omega/t}.
 \label{eq:sign}
\end{equation}
All residue signs below use this convention. Write
\begin{equation}
 L_0=\frac{L_++L_-}{2}.
 \label{eq:linear-median}
\end{equation}
For our meromorphic logarithmic tail this is the principal-value Laplace integral. We call it the \emph{linear median}. An arithmetic average is not, in general, an algebra homomorphism. In particular, inversion does not commute with \eqref{eq:linear-median}. No general nonlinear median-summation construction is being defined by this formula.

\subsection{Three different loci}
A Borel Stokes direction points toward a Borel singularity. A boundary of exponential decay is described by $\Re(A/t)=0$. A critical value of an analytic forward map is a branch point of its inverse. In this article these are, respectively, the positive Borel ray for the physical lateral ambiguity, the imaginary $t$-axis for $Q$, and the moving discriminant in the target variable. Calling all three a ``Stokes line'' would obscure the analysis.

For logarithmic or multivalued leading terms, a complex transseries is interpreted on a chosen sheet and a sectorial asymptotic scale. There is no single dominance ordering of all complex exponentials on the whole punctured plane. Our explicit transseries are generated by powers and logarithms together with $Q$ or $Q^{1/2}$ on such charts. More general action sets require their own support and summability hypotheses.

\section{The shifted product and its Borel transform}
\label{sec:borel}

\subsection{Elementary core and all-orders tail}
Define
\begin{align}
 P_a(t)&=-\frac{A_0}{t}+\left(\frac12-a\right)\log t
       +\frac12\log(2\pi)-\log\Gamma(a)+\frac{B_2(a)}4t,
       \label{eq:core}\\
 \wt R_a(t)&=\sum_{n\geq1}c_{2n}(a)t^{2n},
 &c_{2n}(a)&=-\frac{B_{2n}B_{2n+1}(a)}{2n(2n+1)!}.
 \label{eq:tail}
\end{align}
Here the Bernoulli convention is
\begin{equation}
 \frac{ze^{az}}{e^z-1}=\sum_{n\geq0}B_n(a)\frac{z^n}{n!},
 \qquad B_n=B_n(0).
 \label{eq:bernoulli}
\end{equation}
In particular, $c_2=-B_3(a)/72$ and $c_4=B_5(a)/14400$. The series contains no odd powers beyond the linear term in $P_a$.

\begin{proposition}[Meromorphic Borel kernel]
\label{prop:borel}
For real $0<a<1$, the Borel transform of \eqref{eq:tail} is
\begin{align}
 \widehat R_a(\xi)
 &=-\frac2\pi\sum_{k,m\geq1}\frac{\sin(2\pi ka)}k
                  \frac{\xi}{(Akm)^2-\xi^2}\label{eq:fourier-kernel}\\
 &=\frac1\pi\sum_{k,m\geq1}\frac{\sin(2\pi ka)}k
      \left(\frac1{\xi-Akm}+\frac1{\xi+Akm}\right).
 \nonumber
\end{align}
The paired series converges normally on compact subsets avoiding
$A\Z\setminus\{0\}$. Its residues at both $An$ and $-An$ are
\begin{equation}
 r_a(n)=\frac1\pi\sum_{d\mid n}\frac{\sin(2\pi da)}d,
 \qquad n\geq1.
 \label{eq:residues}
\end{equation}
A residue can vanish; the corresponding apparent pole is then removable.
\end{proposition}
\begin{proof}
Use the absolutely convergent Fourier expansions
\begin{align*}
 B_{2n}&=(-1)^{n+1}\frac{2(2n)!}{(2\pi)^{2n}}
                \sum_{m\geq1}m^{-2n},\\
 B_{2n+1}(a)&=(-1)^{n+1}\frac{2(2n+1)!}{(2\pi)^{2n+1}}
                \sum_{k\geq1}\frac{\sin(2\pi ka)}{k^{2n+1}}
\end{align*}
for $n\geq1$. Substitution in \eqref{eq:tail}, followed by division by
$(2n-1)!$, gives the Taylor coefficients of \eqref{eq:fourier-kernel} at
zero. On a fixed compact set, all sufficiently large $km$ contribute
$O(|\xi|/(k^3m^2))$. The remaining terms are finite in number. This proves
normal convergence of the paired expression and permits continuation.
At $\xi=An$, exactly the pairs $km=n$ have a pole, and their residues sum
to \eqref{eq:residues}. The calculation at $-An$ is the same.
\end{proof}

The Fourier formulas are standard consequences of the Hurwitz-zeta
functional equation; see \cite{dlmfzeta}. The pole data in rational-shift
normalizations already appear in \cite{fr}. They are included here to fix
normalizations for the completion and inverse calculations.

\subsection{A kernel valid for holomorphic shift parameters}
The Fourier series in \eqref{eq:fourier-kernel} must \emph{not} be used
termwise for a nonreal $a$: its sine coefficients then grow exponentially
with $k$. The correct holomorphic continuation is as follows.

\begin{proposition}[Parameter-holomorphic kernel]
\label{prop:holomorphic-kernel}
For all complex $a$, the Borel germ of \eqref{eq:tail} continues to
\begin{equation}
 \widehat R_a(\xi)=\frac1\xi\sum_{m\geq1}
 \left\{
 \frac{\sin\big((a-\frac12)\xi/(2\pi m)\big)}
      {\sin\big(\xi/(4\pi m)\big)}-(2a-1)
 \right\}.
 \label{eq:holomorphic-kernel}
\end{equation}
The value at zero is removable. The sum is locally normally convergent
jointly in $a$ and $\xi$ off the same possible pole set. On a compact
$a$-set it has finite exponential type along rays bounded away from the
real pole directions. Consequently its lateral sums and their $a$
derivatives exist for sufficiently small $t$ on the appropriate sectors.
\end{proposition}
\begin{proof}
The odd Bernoulli generating function follows from
\eqref{eq:bernoulli}; substituting the Fourier formula for $B_{2n}$ then
summing the Bernoulli-polynomial generating function gives
\eqref{eq:holomorphic-kernel} near zero. Alternatively its Taylor
coefficients can be checked directly. The summand in braces is
$O_K(|\xi|^2/m^2)$ uniformly for $a$ in a compact set $K$ and bounded
$\xi$, so normal convergence follows after division by $\xi$.

For a ray $\arg\xi=\theta$ with $|\sin\theta|>0$, split the sum at
$m=1+\lfloor|\xi|\rfloor$. For larger $m$, Taylor's theorem bounds the
tail by $O_K(|\xi|)$. For the finite initial portion, the sine quotient
is bounded by a polynomial factor times $\exp(C_K|\xi|/m)$, hence by a
polynomial factor times $\exp(C_K|\xi|)$. The estimates hold on a small
angular neighborhood of the ray. Laplace convergence follows when
$\Re(\e^{\ii\theta}/t)>C_K$, and Cauchy estimates on a slightly larger
$a$-compact justify parameter differentiation. Contour deformation
defines the corresponding small-$t$ lateral functions. This is a local
statement in $a,t$, not a uniform claim for unbounded complex shifts.
\end{proof}

\begin{remark}[Growth for real shifts]
For $0<a<1$, one can use the aggregated paired kernel
$\sum_n r_a(n)((\xi-An)^{-1}+(\xi+An)^{-1})$ and
$|r_a(n)|\leq(1+\log n)/\pi$. Away from fixed neighborhoods of the poles,
it has at most polynomial growth, with logarithmic factors. Subtracting
the local pole parts gives well-defined principal-value and lateral
Laplace integrals for every $t\in\mathbb H_R$. The proof of the exact
identity below provides a second, normally convergent representation of
these integrals.
\end{remark}

\section{An exact Mellin--Borel completion theorem}
\label{sec:completion}

\subsection{Dual products}
Define the holomorphic functions
\begin{align}
 D_a^\pm(t)&=\log(\e^{\pm2\pi\ii a}Q;Q)_\infty
 =-\sum_{k\geq1}\frac{\e^{\pm2\pi\ii ka}Q^k}{k(1-Q^k)},
 \label{eq:dual}\\
 D_a^0(t)&=\frac{D_a^+(t)+D_a^-(t)}2
 =-\sum_{k\geq1}\frac{\cos(2\pi ka)Q^k}{k(1-Q^k)}.
 \label{eq:dualmedian}
\end{align}
For real $a$, these converge normally on $\mathbb H_R$. In a bounded
complex $a$-neighborhood they converge for all sufficiently small $|Q|$.
The initial factor of each product is $1-\e^{\pm2\pi\ii a}Q$, not
$1-\e^{\pm2\pi\ii a}$; this distinction is responsible for the
logarithmic constants in \cref{cor:fr}.

\begin{theorem}[Exact completion of a shifted infinite product]
\label{thm:completion}
For real $0<a<1$ and every $t\in\mathbb H_R$,
\begin{align}
 G_a(t)&=P_a(t)+L_0\wt R_a(t)+D_a^0(t),
 \label{eq:completion0}\\
 &=P_a(t)+L_+\wt R_a(t)+D_a^-(t),
 \label{eq:completionplus}\\
 &=P_a(t)+L_-\wt R_a(t)+D_a^+(t).
 \label{eq:completionminus}
\end{align}
In particular,
\begin{equation}
 L_+\wt R_a-L_-\wt R_a
 =D_a^+-D_a^-=-2\pi\ii\sum_{n\geq1}r_a(n)Q^n.
 \label{eq:jump}
\end{equation}
The functions $G_a$ themselves have no discontinuity between the two
representations: the change in completion cancels the lateral change.
\end{theorem}

\subsection{The scalar principal-value kernel}
\begin{lemma}[Cotangent Mellin kernel]
\label{lem:kernel}
For $\Re x>0$ put
\begin{equation}
 K(x)=-\frac2\pi\PV\int_0^\infty
             \e^{-xs}\frac{s}{1-s^2}\dd s.
 \label{eq:Kpv}
\end{equation}
For $0<c<2$,
\begin{equation}
 K(x)=\frac1{2\pi\ii}\int_{\Re u=c}
       \Gamma(u)\cot\frac{\pi u}{2}\,x^{-u}\dd u.
 \label{eq:Kmb}
\end{equation}
It also has the expression
\begin{equation}
 K(x)=-\frac1\pi\left(\e^{-x}\Ei(x)-\e^x E_1(x)\right),
 \label{eq:Kei}
\end{equation}
with branches obtained from the positive real $x$-axis.
\end{lemma}
\begin{proof}
For real $x>0$, the Mellin transform in $x$ of \eqref{eq:Kpv} is
\[
 -\frac{2\Gamma(u)}\pi\PV\int_0^\infty
           \frac{s^{1-u}}{1-s^2}\dd s
 =\Gamma(u)\cot\frac{\pi u}{2},\qquad 0<\Re u<2.
\]
Indeed, with $v=s^2$, the principal value is
$-\frac\pi2\cot(\pi u/2)$. This follows, for example, by splitting at
$v=1$, applying $v\mapsto1/v$ to the upper part, and using the partial
fraction expansion of the cotangent.

For completeness, the Mellin interchange is a principal-value
interchange rather than Tonelli's theorem. Remove a symmetric interval
$(1-\delta,1+\delta)$ and, on a fixed neighborhood of $1$, subtract the
value at $1$ from the smooth numerator multiplying $(s-1)^{-1}$. The
subtracted quotient is locally integrable uniformly as $\delta\to0$;
the remaining elementary pole integral has its explicit principal
value. At zero and infinity the powers are integrable precisely in the
strip $0<\Re u<2$. The same subtraction works after integrating in
$x$ on compact Mellin substrips. Mellin inversion yields
\eqref{eq:Kmb}. The gamma factor gives exponential decay on vertical
lines. Both sides extend holomorphically to $\Re x>0$, establishing the
stated complex domain. Partial fractions in \eqref{eq:Kpv} give
\eqref{eq:Kei}.
\end{proof}

\subsection{Mellin proof of the theorem}
\begin{proof}[Proof of \cref{thm:completion}]
Start with real $t>0$. Mellin inversion of the exponentials in
\eqref{eq:G}, with $c_0>1$, gives
\begin{equation}
 G_a(t)=-\frac1{2\pi\ii}\int_{\Re s=c_0}
          \Gamma(s)\zeta(s+1)\zeta(s,a)t^{-s}\dd s.
 \label{eq:mellinG}
\end{equation}
Absolute convergence on that line follows from the original double
sum. Shift to $\Re s=-c$, where $1<c<2$. The three crossed locations are
$s=1,0,-1$. The residues of the full integrand with the preceding minus
sign are
\[
 -\frac{A_0}{t},\qquad
 \left(\frac12-a\right)\log t+\frac12\log(2\pi)-\log\Gamma(a),
 \qquad \frac{B_2(a)}4t.
\]
For the middle expression, use $\zeta(0,a)=\frac12-a$,
$\zeta'(0,a)=\log\Gamma(a)-\frac12\log(2\pi)$, and
$\Gamma(s)\zeta(1+s)=s^{-2}+O(1)$. These standard special values and the
Hurwitz functional equation are recorded in \cite{dlmfzeta}.

The contour shift is justified by gamma decay and polynomial bounds
for the zeta functions on fixed vertical strips. With $s=-u$ the
remaining integral is
\begin{equation}
 G_a-P_a=-\frac1{2\pi\ii}\int_{\Re u=c}
          \Gamma(-u)\zeta(1-u)\zeta(-u,a)t^u\dd u.
 \label{eq:shiftedmellin}
\end{equation}
Introduce the absolutely convergent series
\begin{equation}
 C_z(a)=\sum_{k\geq1}\frac{\cos(2\pi ka)}{k^z},\qquad
 S_z(a)=\sum_{k\geq1}\frac{\sin(2\pi ka)}{k^z},\qquad \Re z>1.
 \label{eq:CS}
\end{equation}
The Riemann and Hurwitz functional equations, followed by gamma
reflection, give the exact identity
\begin{equation}
 \Gamma(-u)\zeta(1-u)\zeta(-u,a)
 =\Gamma(u)A^{-u}\zeta(u)
   \left\{C_{u+1}(a)-\cot\frac{\pi u}{2}S_{u+1}(a)\right\}.
 \label{eq:functionaleq}
\end{equation}
To make its sign transparent, the two zeta factors are
\begin{align*}
 \zeta(1-u)&=2(2\pi)^{-u}\Gamma(u)\cos(\pi u/2)\zeta(u),\\
 \zeta(-u,a)&=\frac{2\Gamma(1+u)}{(2\pi)^{1+u}}
  \{-\sin(\pi u/2)C_{u+1}(a)+\cos(\pi u/2)S_{u+1}(a)\},
\end{align*}
and $\Gamma(-u)\Gamma(1+u)=-\pi/\sin(\pi u)$.

The cosine part of \eqref{eq:shiftedmellin} is now
\[
 -\sum_{k,m\geq1}\frac{\cos(2\pi ka)}k\e^{-Akm/t}=D_a^0(t),
\]
by inverse Mellin transformation of $\Gamma(u)$. The sine--cotangent
part is, by \cref{lem:kernel},
\begin{equation}
 \sum_{k,m\geq1}\frac{\sin(2\pi ka)}k K(Akm/t).
 \label{eq:sumK}
\end{equation}
All these Mellin interchanges are absolutely convergent on
$1<c<2$: the double Dirichlet majorant is
$\sum_{k,m}k^{-c-1}m^{-c}<\infty$, and the gamma factor is integrable
on the vertical line. In particular, \eqref{eq:Kmb} bounds
$|K(Akm/t)|$ on compact subsets of $\mathbb H_R$ by a constant times
$(km)^{-c}$.

Expression \eqref{eq:sumK} is exactly the principal-value Laplace
transform of \eqref{eq:fourier-kernel}. Here is an explicit justification
that avoids an unregularized principal-value interchange. For real
$t>0$, first integrate along a fixed ray $\arg\xi=\theta$, with
$0<|\theta|<\pi/2$. With $\omega=Akm$ and $\xi=r\e^{\ii\theta}$,
\[
 |\omega^2-r^2\e^{2\ii\theta}|
 \geq |\sin\theta|(\omega^2+r^2).
\]
The absolute integral of the $(k,m)$-term is therefore at most a
constant depending on $(t,\theta)$ times $k^{-3}m^{-2}$. Termwise
integration is justified. Deforming each single-pole integral to the
positive ray gives
\[
 L_\pm\wt R_a(t)=\sum_{k,m\geq1}\frac{\sin(2\pi ka)}k
       \{K(Akm/t)\mp\ii\e^{-Akm/t}\}.
\]
Both resulting series converge absolutely: use the Mellin bound for
$K$ and exponential decay for the residues. Their arithmetic average
is \eqref{eq:sumK}. This proves \eqref{eq:completion0} on the positive
real axis.
The normally convergent Mellin representations show that both sides
are holomorphic on $\mathbb H_R$, so the identity theorem proves it
throughout that half-plane.

Finally the residue bound in \eqref{eq:residues} makes
$\sum_nr_a(n)Q^n$ normally convergent. Contour deformation gives the
first equality in \eqref{eq:jump}. Expanding \eqref{eq:dual} as a
geometric series gives the second equality. Since
$L_+=L_0+(D_a^+-D_a^-)/2$, the upper-lateral representation pairs with
$D_a^-$, and the lower-lateral one pairs with $D_a^+$. This proves
\eqref{eq:completionplus} and \eqref{eq:completionminus}.
\end{proof}

\subsection{Asymptotic meaning and endpoints}
The exact theorem implies, and is much stronger than, the fixed-order
asymptotic expansion
\begin{equation}
 G_a(t)=P_a(t)+\sum_{n=1}^{N}c_{2n}(a)t^{2n}+O_{a,N,\eps}(t^{2N+2})
 \quad(t\to0\text{ in }S_{\eps,T}).
 \label{eq:fixedorder}
\end{equation}
This follows from the Taylor expansion of the Borel germ and Watson's
Laplace argument; the pole contributions are flat. The same assertion
also follows by moving the Mellin contour past a prescribed finite
number of negative even integers. The coefficients are Gevrey--1 when
indexed by the full power of $t$. Except at special symmetric shifts,
ordinary convergence is not asserted.

Continuity at $a=1$ and the value $a=1/2$ give
\begin{align}
 G_1(t)&=-\frac{A_0}{t}-\frac12\log t+\frac12\log(2\pi)+\frac t{24}
         +\log(Q;Q)_\infty,
 \label{eq:eta}\\
 G_{1/2}(t)&=-\frac{A_0}{t}+\frac12\log2-\frac t{48}
         +\log(-Q;Q)_\infty.
 \label{eq:half}
\end{align}
Their Bernoulli tails and Stokes jumps vanish identically, but their
flat completions do not. Formula \eqref{eq:eta} is the classical eta
transformation in this normalization, not a new modular law.

\subsection{The exact shift relation}
For $a>0$ the product itself satisfies
\begin{equation}
 G_{a+1}(t)-G_a(t)=-\log(1-\e^{-at}).
 \label{eq:Gshift}
\end{equation}
Moreover,
\begin{align}
 P_{a+1}-P_a&=-\log(at)+\frac{at}{2},\label{eq:Pshift}\\
 \wt R_{a+1}-\wt R_a
 &=-\log(1-\e^{-at})+\log(at)-\frac{at}{2}
 =-\sum_{n\geq1}\frac{B_{2n}(at)^{2n}}{2n(2n)!}.
 \label{eq:Rshift}
\end{align}
The last series is convergent for $|at|<2\pi$ with the specified local
branches. Also $D_{a+1}^\pm=D_a^\pm$. These identities extend the
completion formula to real positive shifts and, using
\cref{prop:holomorphic-kernel}, to a complex neighborhood of any such
shift for sufficiently small sectorial $t$. For the $q$-gamma fold,
one only needs a neighborhood compactly contained in $1<\Re a<2$.
The Fourier sine representation alone would incorrectly discard the
convergent correction in \eqref{eq:Rshift} outside the basic interval.

\subsection{Translation to the recent scalar summability conjecture}
\begin{corollary}[Rational-shift summability identity]
\label{cor:fr}
Fix integers $N\geq2$ and $1\leq k<N$. With $q=\e^{2\pi\ii y}$,
$\Im y>0$, define
\[
 f_{k,N}(y)=\log(\e^{2\pi\ii k/N};q)_\infty,\qquad
 g_{k,N}(y)=\log(q^k;q^N)_\infty.
\]
Use the branches fixed by the logarithmic products. Then
\begin{align}
 s_0(\wt g_{k,N})(y)
 &=g_{k,N}(y)-f_{k,N}\!\left(-\frac1{Ny}\right)
       +\log(1-\e^{2\pi\ii k/N}),&&\Re y>0,\label{eq:fr0}\\
 s_\pi(\wt g_{k,N})(y)
 &=g_{k,N}(y)-f_{N-k,N}\!\left(-\frac1{Ny}\right)
       +\log(1-\e^{-2\pi\ii k/N}),&&\Re y<0.\label{eq:frpi}
\end{align}
Here $s_0,s_\pi$ are Borel sums in the $y$-coordinate, including the
elementary terms of the expansion. These are the formulas of
Conjecture~5 in \cite[version 2, equations (60a)--(60b)]{fr}.
\end{corollary}
\begin{proof}
Set $a=k/N$ and $t=-2\pi\ii Ny$. Then $G_a(t)=g_{k,N}(y)$ and
$Q=\e^{-2\pi\ii/(Ny)}$. Removing the zero-th factor in $f$ gives
\[
 D_a^+(t)=f_{k,N}(-1/(Ny))-\log(1-\e^{2\pi\ii k/N}),
\]
and similarly for $D_a^-$. A positive real $y$-Borel ray becomes a
negative imaginary $t$-Borel ray; deforming it toward the positive
$t$-ray approaches the poles from below. Thus $s_0$ corresponds to
$P_a+L_-\wt R_a$ in the overlap. The negative real $y$-Borel ray
corresponds to the upper lateral $t$-sum. Equations
\eqref{eq:completionminus} and \eqref{eq:completionplus} prove the
identities in the overlap with $\Im y>0$, and analytic continuation
within the stated quadrants completes the proof.
\end{proof}

The assertion concerns the nondegenerate shifts $1\leq k<N$. At
$k=N$ the separate terms involving $\log(1-1)$ are not finite; the
regularized combination is $D_1$, and \eqref{eq:eta} supplies the
endpoint formula. One should not subtract two undefined logarithms.

\section{Reflection kernels and cancellation of exponential actions}
\label{sec:reflection}

\subsection{A necessary-and-sufficient criterion}
Let $S\subset(0,1)$ be a finite set of distinct real shifts. Enlarge it,
if necessary, so that $a\in S$ implies $1-a\in S$, assigning weight zero
to any added point. Let $w_a\in\C$ and put
\begin{equation}
 G_w=\sum_{a\in S}w_aG_a,
 \qquad P_w=\sum_aw_aP_a,
 \qquad \wt R_w=\sum_aw_a\wt R_a.
 \label{eq:weighted}
\end{equation}
Weights need not be rational, multiplicative, or character values.

\begin{theorem}[The two reflection kernels]
\label{thm:reflection}
With the preceding notation:
\begin{equation}
 G_w=P_w+L_0\wt R_w
 \quad\Longleftrightarrow\quad
 w_{1-a}=-w_a\quad\text{for every }a\in S;
 \label{eq:oddcriterion}
\end{equation}
whereas
\begin{equation}
 L_+\wt R_w=L_-\wt R_w
 \quad\Longleftrightarrow\quad
 w_{1-a}=w_a\quad\text{for every }a\in S.
 \label{eq:evencriterion}
\end{equation}
In particular, the coefficient at $a=1/2$ must be zero in
\eqref{eq:oddcriterion}, but is unrestricted in \eqref{eq:evencriterion}.
For nonzero weights the two properties cannot hold simultaneously.
\end{theorem}
\begin{proof}
By \cref{thm:completion}, the error in median reconstruction is
\begin{equation}
 D_w^0=-\sum_{n\geq1}C_n(w)Q^n,
 \qquad C_n(w)=\sum_{d\mid n}\frac{c_d(w)}d,
 \quad c_d(w)=\sum_aw_a\cos(2\pi da).
 \label{eq:cosdivisor}
\end{equation}
It vanishes identically if and only if all $C_n(w)$ vanish. M\"obius
inversion, or induction over divisors, shows this is equivalent to
$c_d(w)=0$ for every positive integer $d$.

For each distinct nonzero phase $z$ among $\e^{2\pi\ii a}$ and its
inverse, collect the coefficients in
\[
 c_d(w)=\frac12\sum_aw_a(z_a^d+z_a^{-d}).
\]
If $z_1,\ldots,z_p$ are the distinct phases, the equations for
$d=1,\ldots,p$ form an invertible Vandermonde system, since every
$z_j\ne0$. Hence the coefficient of each phase is zero. For a
non-self-reflected pair this says $w_a+w_{1-a}=0$; at $a=1/2$ it says
$w_{1/2}=0$. The converse follows immediately from the cosine symmetry.

Similarly, the jump is
\begin{equation}
 L_+\wt R_w-L_-\wt R_w
 =-2\ii\sum_{n\geq1}S_n(w)Q^n,
 \qquad S_n(w)=\sum_{d\mid n}\frac{s_d(w)}d,
 \quad s_d(w)=\sum_aw_a\sin(2\pi da).
 \label{eq:sindivisor}
\end{equation}
The same divisor inversion and Vandermonde argument now gives
$w_a-w_{1-a}=0$. The self-reflected sine is identically zero. Over
$\C$, simultaneous symmetry and antisymmetry force every weight to
vanish.
\end{proof}

\begin{corollary}[Character-weighted median reconstruction]
\label{cor:characters}
For a nonzero Dirichlet character $\chi$ modulo $N\geq2$, the sum
$\sum_{k=1}^{N-1}\chi(k)G_{k/N}(t)$ is exactly reconstructed by the
linear median of its full power--logarithmic expansion if and only if
$\chi(-1)=-1$. Its perturbative Stokes jump vanishes if and only if
$\chi(-1)=1$.
\end{corollary}
\begin{proof}
Use $\chi(N-k)=\chi(-1)\chi(k)$ and
\cref{thm:reflection}. A character is not identically zero on the
reduced residue classes, so the two alternatives are distinct.
\end{proof}
This proves the corresponding $g$-family median-reconstruction
assertion in \cite[Conjecture 1]{fr}, in its logarithmic-product
normalization. Primitivity is not needed for this reconstruction
criterion. This conclusion does not identify every number-theoretic
property required by the separate definition of modular resurgence.

\subsection{Finite-phase bounds for delayed actions}
\begin{proposition}[How long can the first action be delayed?]
\label{prop:delay}
Let $p$ be the number of distinct phases in
$\{\e^{2\pi\ii a},\e^{-2\pi\ii a}:a\in S\}$.
If $D_w^0\not\equiv0$, its first nonzero coefficient is at $Q^r$ for
some $r\leq p$. If the jump is nonzero, its first nonzero coefficient
is at $Q^r$ for some $r\leq p-1$.
\end{proposition}
\begin{proof}
The divisor transform is lower triangular with diagonal coefficient
$1/n$: vanishing of its first $M$ outputs is equivalent to vanishing
of the first $M$ input moments. If the first $p$ cosine moments vanish,
the Vandermonde argument in \cref{thm:reflection} annihilates all phase
coefficients, so the completion is zero. For the sine sequence its
zeroth moment is already zero. Vanishing of moments $1$ through $p-1$
therefore supplies the $p$ consecutive equations $0$ through $p-1$,
again forcing all coefficients to vanish.
\end{proof}

The bound is a finite-data statement, not a bound for arbitrary
infinite combinations. A single shift $a=1/4$ illustrates why even a
one-orbit completion can lose its first action:
\begin{equation}
 C_1=0,\qquad C_2=-\frac12,
 \qquad S_1=1.
 \label{eq:quarter}
\end{equation}
Thus the median reconstruction error begins at $Q^2$ while the Stokes
jump begins at $Q$. At $a=1/2$, the converse contrast is sharper: the
jump is zero but $D_{1/2}^0=Q+O(Q^2)$ is nonzero.

\subsection{What information is invisible?}
Reflection gives a direct decomposition of weights into symmetric and
antisymmetric parts. The tail $\wt R_w$ sees only the antisymmetric
part, because $B_{2n+1}(1-a)=-B_{2n+1}(a)$. The median completion sees
only the symmetric part. The finite elementary core $P_w$ can still
contain information about both parts. Consequently it would be too
strong to say that the entire asymptotic expansion never records any
symmetric data.

A separate example shows that a genuinely nonzero flat function can
have an entirely zero asymptotic expansion. Let $E(t)=G_1(t)$ and
\begin{equation}
 (c_0,c_1,c_2,c_3,c_4)=(2,-9,14,-9,2),\qquad
 Z(t)=\sum_{j=0}^4c_jE(2^jt).
 \label{eq:eta-flat}
\end{equation}
The polynomial
\[
 \sum_{j=0}^4c_jz^j=(z-1)^2(z-2)(2z-1)
\]
implies $\sum c_j2^{-j}=\sum c_j2^j=\sum c_j=\sum jc_j=0$.
All terms of the elementary core in \eqref{eq:eta} cancel. Therefore
\begin{equation}
 Z(t)=\sum_{j=0}^4c_j
       \log(\e^{-A/(2^jt)};\e^{-A/(2^jt)})_\infty
 \sim-2\e^{-A/(16t)}\qquad(t\to0^+).
 \label{eq:eta-flat-asymptotic}
\end{equation}
There is no power-series Borel tail or associated jump from which this
nonzero function could be recovered. This is a classical modular
mechanism presented as an explicit warning for inverse reconstruction,
not as a new eta transformation.

\section{Reversion on formal and analytic charts}
\label{sec:inverse}

\subsection{A precise formal existence mechanism}
A statement about ``general transseries'' needs a specified completion
and a specified substitution operation. The following formulation
isolates what a formal inverse calculation actually uses.

\begin{lemma}[Filtered fixed-point reversion]
\label{lem:filtered}
Let $M=F^1M$ be a complete, separated filtered abelian group with
$F^{n+1}M\subset F^nM$. Suppose a well-defined map $T:M\to M$ satisfies
\begin{equation}
 u-v\in F^nM\quad\Longrightarrow\quad
 T(u)-T(v)\in F^{n+1}M\qquad(n\geq1).
 \label{eq:filteredcontraction}
\end{equation}
Then $u=T(u)$ has exactly one solution in $M$. Iteration from any
initial point determines the solution modulo $F^nM$ after finitely
many steps. The statement applies to a normalized inverse equation
whenever its transseries substitutions are defined and obey
\eqref{eq:filteredcontraction}.
\end{lemma}
\begin{proof}
Start with $u_0\in M$, $u_{n+1}=T(u_n)$. Since
$u_1-u_0\in F^1M$, induction gives
$u_{n+1}-u_n\in F^{n+1}M$. Thus $(u_n)$ is Cauchy and converges by
completeness. Condition \eqref{eq:filteredcontraction} gives continuity
of $T$, so the limit is a fixed point. If $u,v$ are fixed points, their
difference belongs successively to every $F^nM$, and is zero by
separatedness.
\end{proof}

For a chosen leading inverse $h(Y)$, one writes $x=h(Y)+u$ and solves
for $u$. If the derivative of the normalized leading equation is a
unit, Taylor expansion often yields exactly this contraction.
A complete $Q$-adic algebra is an immediate example. For several
exponential actions, a positive action filtration can be used only
when its sums are locally finite and the required substitutions are
continuous. Countably many actions accumulating at zero do not
satisfy these hypotheses merely because each exponential is small.
No unrestricted Hahn-field closure theorem is hidden in
\cref{lem:filtered}.

\subsection{Analytic inverse transport with a convergence certificate}
\begin{theorem}[All-orders transport of an analytic completion]
\label{thm:inverse}
Let $F$ be univalent on a neighborhood $U$, let $h=F^{-1}$ on $V=F(U)$,
and let $E$ be holomorphic on $U$. Fix a target $Y$ and a closed disk
$\overline{D(Y,r)}\subset V$. Put $e(z)=E(h(z))$ and suppose
\begin{equation}
 |\lambda|\sup_{|z-Y|=r}|e(z)|<r.
 \label{eq:rouche}
\end{equation}
There is a unique solution in $h(D(Y,r))$ of
$F(x)+\lambda E(x)=Y$, counted with multiplicity. It is simple,
analytic in $\lambda$ throughout the disk defined by
\eqref{eq:rouche}, and has the convergent expansion
\begin{equation}
 x(Y,\lambda)=h(Y)+\sum_{n\geq1}\frac{(-\lambda)^n}{n!}
 \partial_Y^{n-1}\!\left[h'(Y)\,E(h(Y))^n\right].
 \label{eq:inverseall}
\end{equation}
If $|\lambda|\sup|e|<r$ holds uniformly over compact target sets, the
convergence is locally uniform in the target as well.
\end{theorem}
\begin{proof}
In the coordinate $z=F(x)$ the equation is $z+\lambda e(z)=Y$.
Rouch\'e's theorem gives exactly one zero in $D(Y,r)$, counted with
multiplicity. It must be simple, since a multiple zero would have
multiplicity at least two. The implicit-function theorem then gives
analytic dependence throughout the parameter disk.

For small $\lambda$, the residue formula for $h$ evaluated at the
unique zero is
\[
 x(Y,\lambda)=\frac1{2\pi\ii}\int_{|z-Y|=r}
 h(z)\frac{1+\lambda e'(z)}{z+\lambda e(z)-Y}\dd z.
\]
Subtract its value at $\lambda=0$. Integrate the derivative of
$\log(1+\lambda e(z)/(z-Y))$ by parts around the circle, and expand
that logarithm by the absolutely convergent series guaranteed by
\eqref{eq:rouche}. Cauchy's differentiation formula gives
\eqref{eq:inverseall}. Analytic continuation in the parameter disk
and compact uniformity complete the proof.
\end{proof}

This is Lagrange--B\"urmann inversion applied in a forward-value
coordinate. It is included rather than claimed as new; the repository
already develops its multi-parameter form \cite{repotransport}. A
simple quantitative tail bound is also useful. Write
$m_e=\sup_{|z-Y|=r}|e(z)|$ and
$M_h=\sup_{|z-Y|=r}|h'(z)|$. Cauchy's estimate in
\eqref{eq:inverseall} gives
\begin{equation}
 \left|[\lambda^n]x\right|\leq\frac{M_h r}{n}
                       \left(\frac{m_e}{r}\right)^n.
 \label{eq:coefbound}
\end{equation}
For $\chi=|\lambda|m_e/r<1$, the tail above degree $N$ is at most
\begin{equation}
 \frac{M_hr\,\chi^{N+1}}{(N+1)(1-\chi)}.
 \label{eq:tailbound}
\end{equation}
A fixed target disk and derivative control are essential: a bare
pointwise assertion $E(h(Y))\to0$ does not by itself justify a uniform
inverse expansion.

The first two terms are
\begin{equation}
 x=h-\lambda\frac{E(h)}{F'(h)}
 +\lambda^2\left\{\frac{E(h)E'(h)}{F'(h)^2}
                 -\frac{F''(h)E(h)^2}{2F'(h)^3}\right\}
 +O(\lambda^3).
 \label{eq:inverse2}
\end{equation}
The derivative in the second term is one reason why inversion creates
mixed exponential sectors even from a simple additive jump.

\subsection{Exact Stokes transport and its cocycle}
Suppose two lateral forward charts satisfy $F_+=F_-+J$. Let
$h_-=F_-^{-1}$ and set
\begin{equation}
 \Phi(Y)=Y+J(h_-(Y)).
 \label{eq:Phi}
\end{equation}
Where the maps are defined and univalent,
\begin{equation}
 F_+=\Phi\circ F_-,\qquad
 h_+=h_-\circ\Phi^{-1}.
 \label{eq:cocycle}
\end{equation}
This is an exact relation, not just its linearized jump. Formula
\eqref{eq:inverseall} with $E=J$ supplies every inverse jump sector.
For three charts $F_i$, the maps
$\Phi_{ji}=F_j\circ F_i^{-1}$ satisfy
$\Phi_{ki}=\Phi_{kj}\circ\Phi_{ji}$ wherever all compositions make
sense. The inverse charts consequently transform contravariantly.
This elementary cocycle bookkeeping prevents inconsistent branch
choices during analytic continuation.

\subsection{Why arithmetic median and inversion do not commute}
Let $F_\pm=F_0\pm\lambda J/2$, with inverses $h_\pm$ and $h_0=F_0^{-1}$.
Averaging \eqref{eq:inverseall} eliminates odd powers of $\lambda$ but
not the even powers:
\begin{equation}
 \frac{h_++h_-}{2}-h_0
 =\frac{\lambda^2}{8}\partial_Y
        \left[h_0'(Y)J(h_0(Y))^2\right]+O(\lambda^4).
 \label{eq:medianbias}
\end{equation}
Equivalently the coefficient of $\lambda^2$ is
\[
 \frac{J(h_0)J'(h_0)}{4F_0'(h_0)^2}
 -\frac{F_0''(h_0)J(h_0)^2}{8F_0'(h_0)^3}.
\]
The remainder is uniform on smaller target disks under the analytic
bounds of \cref{thm:inverse}. Describing it merely as $O(J^4)$ without
a norm and derivative domain can be misleading. If a nonlinear,
algebra-compatible summation scheme is available, inverse
compatibility follows from uniqueness within that scheme. It does
not follow from arithmetic averaging.

\subsection{Formal coefficients versus selected analytic functions}
General resurgent closure under inversion is established under
appropriate analytic-continuation and convolution bounds
\cite{sauzin,sauzinintro}. In our applications it is important to
separate three convergences: the Borel germ converges near zero; the
exponential-amplitude expansion of \cref{thm:inverse} converges on its
specified parameter disk; and the power-series expansion inside an
exponential sector may be only Gevrey asymptotic. None of the first
two facts implies convergence of the third.

An exact completion can also contain a jump-free part not encoded in
the ordinary formal series. The examples in \cref{sec:reflection}
therefore rule out any unrestricted rule that reconstructs an exact
inverse from the forward power series and the forward jump alone.
Additional normalization, modular data, or a specified analytic
function is needed.

\section{Regular inverse transseries at the \texorpdfstring{$q=1$}{q=1} cusp}
\label{sec:cusp}

\subsection{The logarithmic target and its exact core}
Let $0<a<1$ be fixed and set
\begin{equation}
 x=1/t,\qquad Y=-G_a(1/x).
 \label{eq:target}
\end{equation}
For large positive $x$, $Y$ tends to positive infinity. Introduce the
median carrier and its completion:
\begin{align}
 H_a(x)&=-P_a(1/x)-L_0\wt R_a(1/x),\label{eq:H}\\
 E_a(x)&=-D_a^0(1/x)=\sum_{n\geq1}C_n(a)\e^{-Anx},
 &C_n(a)&=\sum_{d\mid n}\frac{\cos(2\pi da)}d.
 \label{eq:Ecusp}
\end{align}
Thus the exact forward map is $H_a+E_a$. Its algebraic expansion begins
\begin{equation}
 H_a(x)\sim A_0x+b\log x+c+\frac d{x}
                       -\sum_{n\geq1}c_{2n}(a)x^{-2n},
 \label{eq:Hasy}
\end{equation}
where
\begin{equation}
 b=\frac12-a,\qquad
 c=\log\Gamma(a)-\frac12\log(2\pi),\qquad d=-\frac{B_2(a)}4.
 \label{eq:bcd}
\end{equation}
On every sufficiently far-out proper closed subsector of $\Re x>0$,
$H_a'=A_0+o(1)$ uniformly on fixed-radius disks. The inverse theorem therefore supplies a coherent
large real inverse branch $h_a$ and sectorial continuations of that
branch on smaller angular domains. Statements below always refer to
these charts, not to all solutions of the transcendental equation.

A convenient exact leading core $\ell(Y)$ solves
\begin{equation}
 A_0\ell+b\log\ell+c=Y.
 \label{eq:lincore}
\end{equation}
For $b\ne0$, it can be written locally as
\begin{equation}
 \ell(Y)=\frac b{A_0}
 W_j\!\left(\frac{A_0}{b}\exp\frac{Y-c}{b}\right),
 \label{eq:Lambert}
\end{equation}
with the branch chosen to satisfy \eqref{eq:lincore}, not merely its
exponentiated version. On the large positive real branch, $j=0$ for
$b>0$ and $j=-1$ for $b<0$. If $b=0$, use $\ell=(Y-c)/A_0$.

Writing $x=\ell+u$ and $\eps=1/\ell$ gives the normalized carrier
inverse equation
\begin{equation}
 A_0u+b\log(1+\eps u)+\frac{d\eps}{1+\eps u}
 -L_0\wt R_a\!\left(\frac{\eps}{1+\eps u}\right)=0.
 \label{eq:exactcorenormalization}
\end{equation}
Its formal version has a unique ordinary formal power series
$u(\eps)=O(\eps)$ by \cref{lem:filtered}, since $A_0\ne0$.
With $L_+$ or $L_-$ in place of $L_0$, the summable implicit-function
and composition results in \cite{sauzin} apply to the meromorphic,
finite-exponential-type Borel kernel established above. They provide
the lateral analytic realizations of this ordinary resurgent
correction. The median-carrier inverse itself is not asserted to be
the arithmetic average of those lateral inverses.

This normalization is useful precisely because it does not treat
$\log Y$ as a constant coefficient while performing a Borel transform.
The logarithm is absorbed into the exact core first. Re-expanding
$\ell(Y)$ afterward gives the usual power--logarithmic transseries.

\subsection{Algebraic reversion and fixed-order error control}
Put $L=\log(Y/A_0)$ and $s=bL+c$. Elementary substitution gives
\begin{equation}
 h_a(Y)=\frac{Y-s}{A_0}
       +\frac{bs-A_0d}{A_0Y}
       +O\!\left(\frac{(\log Y)^2}{Y^2}\right).
 \label{eq:halgebraic}
\end{equation}
Every further fixed algebraic order can be obtained by substituting a
power--logarithmic ansatz into \eqref{eq:Hasy} and cancelling the next
residual. The coefficient of the new unknown is $A_0$, so the recursion
is triangular. On the chosen chart, a residual $\rho_N(Y)$ yields an
inverse error $O(\rho_N(Y))$, because $H_a'$ stays bounded away from
zero. This proves the fixed-order expansion, rather than merely
producing a formal list of coefficients. Derivative estimates follow
on smaller complex disks by Cauchy's formula.

\subsection{A convergent expansion in the inverse exponential scale}
\begin{theorem}[One- and two-action inverse coefficients]
\label{thm:cusp}
Let $h=h_a(Y)$ and $z=\e^{-Ah}$. Write
$f=H_a'(h)$ and $f_2=H_a''(h)$. For the large real inverse branch, and
uniformly on its smaller sectorial charts, the inverse $x_*(Y)$ of the
exact forward map $H_a+E_a$ has a locally convergent expansion
\begin{equation}
 x_*(Y)=h+\sum_{n\geq1}\Delta_n(h)z^n.
 \label{eq:cuspsectors}
\end{equation}
The first two coefficients are
\begin{align}
 \Delta_1&=-\frac{C_1(a)}f,\label{eq:Delta1}\\
 \Delta_2&=-\frac{C_2(a)}f-\frac{A C_1(a)^2}{f^2}
                         -\frac{f_2C_1(a)^2}{2f^3},
 \label{eq:Delta2}
\end{align}
where $C_1(a)=\cos(2\pi a)$ and
$C_2(a)=\cos(2\pi a)+\frac12\cos(4\pi a)$.
For $h$ sufficiently large in these charts, a common positive radius
in the auxiliary variable $z$ can be chosen.
\end{theorem}
\begin{proof}
Treat $z$ as independent and write $x=h+\delta$. The exact equation is
\begin{equation}
 H_a(h+\delta)-H_a(h)
 +\sum_{n\geq1}C_n(a)z^n\e^{-An\delta}=0.
 \label{eq:independentz}
\end{equation}
On a fixed small $\delta$-disk, the last series converges normally for
$|z|$ sufficiently small, because $C_n(a)=O(\log n)$. The derivative
with respect to $\delta$ at $(\delta,z)=(0,0)$ is $f$, bounded away
from zero for large $h$. The analytic implicit-function theorem gives
\eqref{eq:cuspsectors}. Uniform bounds on a slightly larger disk, or
Rouch\'e's theorem as in \cref{thm:inverse}, give a common $z$-radius.
Equating coefficients of $z$ and $z^2$ in
\eqref{eq:independentz} gives
\[
 f\Delta_1+C_1=0,\qquad
 f\Delta_2+\tfrac12 f_2\Delta_1^2-A C_1\Delta_1+C_2=0,
\]
which proves the formulas.
\end{proof}

The term $-AC_1^2/f^2$ records the change of the exponential argument
under inversion. Keeping only $E_a(h)$ and never expanding
$E_a(h+\delta)$ would miss it. If the first nonzero completion
coefficient is $C_r$, then $\Delta_1=\cdots=\Delta_{r-1}=0$ and
$\Delta_r=-C_r/f$. This observation applies to finite weighted products
as well, provided their leading forward slope is nonzero; their action
scale is then computed using that slope, not automatically $A_0$.

\subsection{Action 24 and its power--logarithmic dressing}
Since $A/A_0=24$, \eqref{eq:halgebraic} gives
\begin{equation}
 \e^{-Ah_a(Y)}=
 \e^{-24Y}(Y/A_0)^{24b}\e^{24c}
 \left\{1+O\!\left(\frac{\log Y}{Y}\right)\right\}.
 \label{eq:action24}
\end{equation}
For $C_1(a)\ne0$, therefore,
\begin{equation}
 x_*(Y)-h_a(Y)
 =-\frac{\cos(2\pi a)}{A_0}\,
 \e^{-24Y}(Y/A_0)^{24b}\e^{24c}
 \left\{1+O\!\left(\frac{\log Y}{Y}\right)\right\}.
 \label{eq:inverseleading}
\end{equation}
The first action is $24r$ when the first surviving completion harmonic
is $r$. In particular $a=1/4$ has first completion action $48$, whereas
its lateral inverse Stokes transition has first action $24$. The
action ratio itself is an elementary consequence of the chosen
normalization; the completion coefficients determine which harmonics
actually occur.

For the original variable $t=1/x$, the first correction is
\begin{equation}
 t_*(Y)=\frac1h+\frac{C_1(a)}{f h^2}\e^{-Ah}
              +O\!\left(\frac{\e^{-2Ah}}{h^2}\right).
 \label{eq:tinverse}
\end{equation}
A final substitution $q=\e^{-t}$ gives the inverse in $q$. All
remainders here refer to the exact median-carrier inverse $h$;
replacing $h$ by a short algebraic truncation creates an additional,
usually much larger, algebraic error.

\subsection{The lateral inverse transition}
Define $H_{a,\pm}=-P_a(1/x)-L_\pm\wt R_a(1/x)$. By
\eqref{eq:jump},
\begin{equation}
 H_{a,+}(x)-H_{a,-}(x)
 =2\ii\sum_{n\geq1}\left(\sum_{d\mid n}\frac{\sin(2\pi da)}d\right)
          \e^{-Anx}.
 \label{eq:Hjump}
\end{equation}
Let $h_{a,\pm}$ be the corresponding large inverse branches. The exact
cocycle \eqref{eq:cocycle} gives their finite transition. Its leading
term is
\begin{equation}
 h_{a,+}-h_{a,-}
 =-\frac{2\ii\sin(2\pi a)}{H_a'(h_a)}\e^{-Ah_a}
       +O(\e^{-2Ah_a}).
 \label{eq:hjump}
\end{equation}
These are inverse branches of two \emph{carriers}, not two different
values of the exact product inverse. Pairing $H_{a,+}$ with $-D_a^-$
or $H_{a,-}$ with $-D_a^+$ reconstructs the same exact forward map and
therefore the same chosen exact inverse. This is the operational
meaning of Stokes consistency for the completed transseries.


\subsection{An exact quadratic carrier with no Stokes jump}
At $a=1/2$, both $\wt R_a$ and its Stokes jump vanish. Nevertheless the
exact inverse differs from the inverse of its elementary core. Here
\[
 H_{1/2}(x)=A_0x-\tfrac12\log2+\frac1{48x},\qquad
 E_{1/2}(x)=-\log(-\e^{-Ax};\e^{-Ax})_\infty.
\]
Writing $Z=Y+\frac12\log2$, the large carrier inverse is explicitly
\begin{equation}
 h_{1/2}(Y)=\frac{Z+\sqrt{Z^2-A_0/12}}{2A_0}.
 \label{eq:quadraticcarrier}
\end{equation}
Its completed inverse begins
\begin{equation}
 x_*(Y)=h_{1/2}(Y)
 +\frac{\e^{-Ah_{1/2}(Y)}}{A_0-1/(48h_{1/2}(Y)^2)}
 +O(\e^{-2Ah_{1/2}(Y)}).
 \label{eq:halfcompletedinverse}
\end{equation}
There is no perturbative Stokes jump whose removal could supply this
term. The exact dual product supplies it. The square root in
\eqref{eq:quadraticcarrier} is a separate algebraic branch choice,
again not a Borel Stokes phenomenon.

\subsection{Cancellation of the linear core raises the transseries height}
The action $24$ is not universal even within shifted products. Fix
$0<a<1/2$ and consider the antisymmetric quotient
\begin{equation}
 F^{\mathrm{odd}}_a(x)=-G_a(1/x)+G_{1-a}(1/x).
 \label{eq:oddforward}
\end{equation}
Its median completion vanishes by \cref{thm:reflection}, but its
$t^{-1}$ core also cancels. Put
\[
 B=1-2a>0,\qquad c_o=\log\Gamma(a)-\log\Gamma(1-a).
\]
Then the exact function is
\begin{equation}
 F^{\mathrm{odd}}_a(x)=B\log x+c_o-2L_0\wt R_a(1/x).
 \label{eq:oddcore}
\end{equation}
For real $Y\to+\infty$, its inverse satisfies
\begin{equation}
 h_o(Y)=\ell_o(Y)+O(\ell_o(Y)^{-1}),\qquad
 \ell_o(Y)=\exp\!\left(\frac{Y-c_o}{B}\right).
 \label{eq:oddinverse}
\end{equation}
Indeed the carrier error is $O(x^{-2})$ and the leading derivative is
$B/x$, so the required correction is $O(x^{-1})$. A normalized
implicit equation in $x/\ell_o$ proves every further fixed order.
Complex statements hold on horizontal target strips that map under
$\ell_o$ into a proper right-half-plane sector, with the logarithmic
sheet specified.

The two lateral forward carriers differ by
$4\ii\sin(2\pi a)\e^{-Ax}+O(\e^{-2Ax})$. Their inverse transition
therefore has the leading form
\begin{equation}
 h_{o,+}-h_{o,-}
 =-\frac{4\ii\sin(2\pi a)}B\,
    \ell_o(Y)\exp(-A\ell_o(Y))\big\{1+O(\ell_o(Y)^{-1})\big\}.
 \label{eq:doubleexponential}
\end{equation}
A single exponential in $1/t$ has become a doubly exponential scale
in $Y$. This is a concrete reason to formulate general reversion by
exact cores and sectorial charts rather than by one fixed set of
exponentials in the original target variable. The exact inverse is
still the inverse of the median forward function, not the arithmetic
average of the two lateral inverses.

\section{Exact completion of the \texorpdfstring{$q$}{q}-gamma function}
\label{sec:qgamma}

\subsection{The normalized function and its expansion}
For $q=\e^{-t}$, define
\begin{equation}
 \Gamma_q(a)=\frac{(q;q)_\infty}{(q^a;q)_\infty}(1-q)^{1-a},
 \qquad
 \cL(a,t)=\log\Gamma_{\e^{-t}}(a).
 \label{eq:qgamma}
\end{equation}
This is the standard $q$-gamma normalization \cite{dlmfqgamma}. On our
local branches,
\begin{equation}
 \cL(a,t)=G_1(t)-G_a(t)+(1-a)\log(1-\e^{-t}).
 \label{eq:qgammalog}
\end{equation}
The $t^{-1}$ and $\log t$ terms cancel. Its formal expansion is
\begin{equation}
 \wt\cL(a,t)=\log\Gamma(a)+\ell_1(a)t
                         +\sum_{n\geq1}\ell_{2n}(a)t^{2n},
 \label{eq:qgammaformal}
\end{equation}
where
\begin{align}
 \ell_1(a)&=-\frac{(a-1)(a-2)}4,\label{eq:ell1}\\
 \ell_{2n}(a)&=\frac{B_{2n}}{2n(2n)!}
          \left\{\frac{B_{2n+1}(a)}{2n+1}+1-a\right\}.
 \label{eq:elln}
\end{align}
For example,
\begin{equation}
 \ell_2(a)=\frac{(a-1)(a-2)(2a+3)}{144}.
 \label{eq:ell2}
\end{equation}
At $a=3$, the exact identity $\Gamma_q(3)=1+q$ checks the first terms:
$\log(1+\e^{-t})=\log2-t/2+t^2/8+O(t^4)$. At $a=1$ and $a=2$ all
corrections vanish, consistently with $\Gamma_q(1)=\Gamma_q(2)=1$.

Let $\cL_0(a,t)$ be the linear-median carrier obtained by applying
$L_0$ to the tail in \eqref{eq:qgammalog}, while retaining the convergent
factor $\log(1-\e^{-t})$ with its specified branch. Equivalently it is
the linear median of \eqref{eq:qgammaformal}. The shift relation in
\cref{sec:completion} fixes this definition near $1<\Re a<2$.

\begin{theorem}[Exact $q$-gamma completion]
\label{thm:qgamma}
In a complex neighborhood of any positive real $a$, for sufficiently
small sectorial $t$,
\begin{equation}
 \cL(a,t)=\cL_0(a,t)+\cE(a,Q(t)),
 \label{eq:qgammacompletion}
\end{equation}
where
\begin{align}
 \cE(a,z)&=\sum_{n\geq1}e_n(a)z^n
 =\sum_{k\geq1}\frac{\big(\cos(2\pi ka)-1\big)z^k}{k(1-z^k)},
 \label{eq:calE}\\
 e_n(a)&=\sum_{d\mid n}\frac{\cos(2\pi da)-1}{d}.
 \label{eq:en}
\end{align}
This is a convergent holomorphic function of $(a,z)$ for $a$ in a
fixed compact neighborhood and $|z|$ sufficiently small. Its first
coefficients are
\begin{equation}
 e(a):=e_1(a)=\cos(2\pi a)-1,\qquad
 e_2(a)=e(a)+\tfrac12\big(\cos(4\pi a)-1\big).
 \label{eq:e12}
\end{equation}
\end{theorem}
\begin{proof}
Subtract \eqref{eq:completion0} for $G_a$ from the endpoint formula for
$G_1$. The additional logarithm in \eqref{eq:qgammalog} is convergent
at $t=0$ after removing $\log t$, so it introduces no lateral
ambiguity. The completion is $D_1^0-D_a^0$, which is exactly
\eqref{eq:calE}. Normal convergence follows from
$|\cos(2\pi ka)|\leq\exp(2\pi k|\Im a|)$ and a choice
$|z|\exp(2\pi\sup|\Im a|)<1$. Uniform differentiation in $a$ follows
on smaller neighborhoods.
\end{proof}

\subsection{Uniformity in the parameter}
Fix $\eps>0$ and a small closed disk $D$ around a positive real
$a_*>1$ with $\overline D\subset\{1<\Re a<2\}$. Shrink $D$, if needed,
so that the products are defined for $t\in S_{\eps,T}$. By the
holomorphic Borel kernel and its shift relation, the expansion
\eqref{eq:qgammaformal} holds uniformly on a slightly larger disk,
with every fixed number of $a$ derivatives. In particular,
\begin{equation}
 \partial_a^j\cL_0(a,t)=\partial_a^j\log\Gamma(a)+O(t)
 \quad (a\in D)
 \label{eq:uniformgamma}
\end{equation}
for each fixed $j$. This can also be obtained from
\eqref{eq:qgammacompletion}: the completion and its fixed derivatives
are flat, and the finite-order expansion follows from the uniformly
holomorphic Borel representation. Cauchy's estimate on a larger
$a$-disk is what justifies differentiating the asymptotic expansion.

\subsection{Real convexity and the actual critical point}
For real $t>0$, termwise differentiation of
\eqref{eq:qgammalog} gives
\begin{align}
 \partial_a\cL(a,t)
  &=-t\sum_{n\geq0}\frac1{\e^{(a+n)t}-1}
                   -\log(1-\e^{-t}),\label{eq:digammaq}\\
 \partial_a^2\cL(a,t)
  &=t^2\sum_{n\geq0}
           \frac{\e^{(a+n)t}}{(\e^{(a+n)t}-1)^2}>0.
 \label{eq:convexityq}
\end{align}
Thus $\cL(\cdot,t)$ is strictly convex on $a>0$. Its derivative tends
to $-\infty$ at $a=0^+$ and to $-\log(1-\e^{-t})>0$ at infinity.
There is exactly one real critical point. Since
$\cL(1,t)=\cL(2,t)=0$, strict convexity places it in $(1,2)$ and makes
it a simple minimum. These conclusions hold for every real $t>0$;
the following asymptotic analysis identifies its exponentially small
displacement from the median carrier as $t\to0$.

\section{The exponentially shifted critical fold}
\label{sec:fold}

\subsection{The carrier critical point and its all-orders recursion}
Let $a_*$ be the unique positive zero of the digamma function $\psi$.
It lies in $(1,2)$ because $\psi(1)=-\gamma<0$,
$\psi(2)=1-\gamma>0$, and
$\psi_1(a)=\sum_{n\geq0}(n+a)^{-2}>0$ for $a>0$.
Write $\kappa_*=\psi_1(a_*)$. Numerical values, used only for
orientation, are
\begin{equation}
 a_*\approx1.4616321449683623,\qquad
 \kappa_*\approx0.9676722454476212.
 \label{eq:astar}
\end{equation}

For sufficiently small $t$ in a proper sector,
\eqref{eq:uniformgamma} and the analytic implicit-function theorem
give a unique carrier critical point $\alpha_0(t)$ near $a_*$.
Define
\begin{equation}
 \partial_a\cL_0(\alpha_0,t)=0,\qquad
 v_0(t)=\cL_0(\alpha_0,t),\qquad
 \kappa(t)=\partial_a^2\cL_0(\alpha_0,t),\qquad
 \kappa_3(t)=\partial_a^3\cL_0(\alpha_0,t).
 \label{eq:carriercritical}
\end{equation}
The function $\kappa(t)$ stays bounded away from zero and tends to
$\kappa_*$. The critical point and critical value have fixed-order
asymptotic expansions obtained by implicit substitution in
\eqref{eq:qgammaformal}. The first terms are
\begin{align}
 \alpha_0(t)&=a_*+a_1t+O(t^2),
 &a_1&=\frac{2a_*-3}{4\kappa_*},\label{eq:alpha0asy}\\
 v_0(t)&=\log\Gamma(a_*)+\ell_1(a_*)t
 +\left\{\ell_2(a_*)-\frac{\ell_1'(a_*)^2}{2\kappa_*}\right\}t^2
 +O(t^3).
 \label{eq:v0asy}
\end{align}
For example, if $\alpha_0=a_*+a_1t+a_2t^2+\cdots$, then
\begin{equation}
 a_2=-\frac{\frac12\psi_2(a_*)a_1^2+
                   \ell_1''(a_*)a_1+\ell_2'(a_*)}{\kappa_*}.
 \label{eq:a2}
\end{equation}
At every subsequent order the new coefficient has multiplier
$\kappa_*\ne0$, giving an all-orders triangular recursion. The
uniform derivative bound transports each finite residual to an
actual critical-point error, so this is a fixed-order asymptotic
construction and not an assertion that the full $t$-series converges.

\subsection{Promoting the flat scale to an independent variable}
The decisive device is to regard $Q$ as independent before substituting
its exponentially small value. Define
\begin{equation}
 \mathscr L(a,t,z)=\cL_0(a,t)+\cE(a,z).
 \label{eq:auxfamily}
\end{equation}
It is holomorphic in $(a,z)$ on a common polydisk for all sufficiently
small sectorial $t$. Its actual $q$-gamma specialization is
$z=Q(t)$. Uniformity in $t$ follows from \eqref{eq:uniformgamma} and
the normally convergent series \eqref{eq:calE}.

\begin{theorem}[Completed critical point and discriminant]
\label{thm:critical}
There are functions $\alpha(t,z)$ and $v(t,z)$, holomorphic in $z$ on
a disk of radius independent of small sectorial $t$, such that
\begin{equation}
 \partial_a\mathscr L(\alpha(t,z),t,z)=0,\qquad
 v(t,z)=\mathscr L(\alpha(t,z),t,z),
 \label{eq:completedcritical}
\end{equation}
with $\alpha(t,0)=\alpha_0(t)$ and $v(t,0)=v_0(t)$. Their expansions
begin
\begin{align}
 \alpha(t,z)&=\alpha_0-\frac{e'(\alpha_0)}{\kappa}z+O(z^2),
 \label{eq:alphashift}\\
 v(t,z)&=v_0+e(\alpha_0)z
 +\left\{e_2(\alpha_0)-\frac{e'(\alpha_0)^2}{2\kappa}\right\}z^2
 +O(z^3).
 \label{eq:vshift}
\end{align}
The remainders are uniform in $t$ on proper sectors. In particular,
$\alpha(t,Q(t))$ and $v(t,Q(t))$ are the actual $q$-gamma critical
point and value near $a_*$. For positive real $t$ the critical point
is the unique minimum described in \eqref{eq:convexityq}.
\end{theorem}
\begin{proof}
The derivative of $\partial_a\mathscr L$ with respect to $a$ at
$(\alpha_0,0)$ is $\kappa$, bounded away from zero. Uniform bounds on
a larger $a$-disk give a common implicit-function radius in $z$.
Equivalently, Rouch\'e's theorem on a small $a$-circle gives a unique
critical point, and the implicit-function theorem makes it
holomorphic. Composition gives $v(t,z)$.

Write $\alpha=\alpha_0+b_1z+O(z^2)$. The critical-point equation gives
$\kappa b_1+e'(\alpha_0)=0$. In the value equation, the coefficient of
$z^2$ is
$\kappa b_1^2/2+e'(\alpha_0)b_1+e_2(\alpha_0)$, proving
\eqref{eq:vshift}. Cauchy estimates in the common $z$-disk give the
uniform remainder statements. Substitution $z=Q(t)$ and
\cref{thm:qgamma} identify the exact critical point and value.
\end{proof}

Since $e'(a)=-2\pi\sin(2\pi a)$, the leading displacement is
\begin{equation}
 \alpha(t,Q)-\alpha_0(t)
 =\frac{2\pi\sin(2\pi\alpha_0(t))}{\kappa(t)}Q+O(Q^2).
 \label{eq:alphashifttrig}
\end{equation}
More importantly, $e(a_*)<0$. The true critical value is lower than
the median critical value by a nonzero multiple of $Q$ for all
sufficiently small positive $t$.

\subsection{Uniform critical-layer inverse transseries}
\begin{theorem}[The action-halving layer]
\label{thm:layer}
Let $\eta$ range over a compact subset of a simply connected domain
avoiding $e(a_*)$, and choose the two local square-root sheets.
For all sufficiently small sectorial $t$, solve
\begin{equation}
 \cL(a,t)=v_0(t)+Q(t)\eta.
 \label{eq:layerlevel}
\end{equation}
The two solutions near $\alpha_0(t)$ satisfy
\begin{align}
 a_\pm(t,\eta)
 ={}&\alpha_0(t)
 \ \pm Q^{1/2}\sqrt{\frac{2(\eta-e(\alpha_0))}{\kappa}}\nonumber\\
 &-Q\left\{\frac{e'(\alpha_0)}{\kappa}
       +\frac{\kappa_3\big(\eta-e(\alpha_0)\big)}{3\kappa^2}\right\}
 +O(Q^{3/2}).
 \label{eq:layerformula}
\end{align}
Here $Q^{1/2}$ means the analytic exponential $\e^{-A/(2t)}$ on the
chosen $t$-sector. The expansion is convergent as a power series in
an auxiliary $s=\sqrt z$ for fixed $(t,\eta)$ in these compact sets,
with a common small radius; substituting $z=Q(t)$ gives the stated
transseries. The remainder is uniform on those compact sets.
\end{theorem}
\begin{proof}
In the auxiliary family \eqref{eq:auxfamily}, put $z=s^2$ and
$a=\alpha_0+su$. After subtracting $v_0+s^2\eta$ and dividing by
$s^2$, the equation extends holomorphically to $s=0$ and becomes
\begin{equation}
 \frac\kappa2u^2+e(\alpha_0)-\eta
 +s\left\{\frac{\kappa_3}{6}u^3+e'(\alpha_0)u\right\}
 +O(s^2)=0.
 \label{eq:scaledfold}
\end{equation}
The two $s=0$ roots
$u_0=\pm\sqrt{2(\eta-e(\alpha_0))/\kappa}$ are uniformly separated
from zero. The implicit-function theorem therefore gives a convergent
series $u=u_0+s u_1+O(s^2)$ uniformly on the stated parameter sets.
The coefficient of $s$ in \eqref{eq:scaledfold} is
\[
 \kappa u_0u_1+\frac{\kappa_3}{6}u_0^3+e'(\alpha_0)u_0=0.
\]
Since $u_0\ne0$,
\[
 u_1=-\frac{e'(\alpha_0)}\kappa
      -\frac{\kappa_3(\eta-e(\alpha_0))}{3\kappa^2}.
\]
Multiplying by $s$ and substituting $s=\e^{-A/(2t)}$ proves
\eqref{eq:layerformula}. Uniform analytic bounds supply the remainder
and convergence assertions.
\end{proof}

The theorem explicitly excludes the moving scaled discriminant.
It must not be described as uniform while $\eta\to e(\alpha_0)$.
The next subsection gives the correct representation at that
coalescence.

\begin{corollary}[Two real inverse values at the carrier minimum]
\label{cor:tworeal}
For all sufficiently small real $t>0$, the equation
$\cL(a,t)=v_0(t)$ has two distinct real solutions near $a_*$. Their
separation is
\begin{equation}
 a_+(t,0)-a_-(t,0)
 =2\sqrt{\frac{2(1-\cos(2\pi\alpha_0(t)))}{\kappa(t)}}
         \e^{-2\pi^2/t}+O(\e^{-6\pi^2/t}).
 \label{eq:realgap}
\end{equation}
Thus a jump-free part of the forward completion is already sufficient
to produce a leading inverse effect at half the forward action.
\end{corollary}
\begin{proof}
At $\eta=0$, the quantity under the square root in
\eqref{eq:layerformula} is positive for small real $t$, because
$\kappa(t)>0$ and $\alpha_0(t)\in(1,2)$. The common $Q$ term cancels
in the difference. Reality follows from the real analytic equations
and the real initial roots in \eqref{eq:scaledfold}. Strict convexity
also identifies these as the two local sides of the unique minimum.
\end{proof}

\subsection{A convergent normal form at the actual branch point}
\begin{theorem}[Uniform completed quadratic normal form]
\label{thm:normalform}
After shrinking the parameter neighborhoods, there is a function
$\Psi(w,t,z)$, holomorphic in $(w,z)$ on common neighborhoods and with
$\Psi(0,t,z)=0$, such that the inverse of the auxiliary family at a
target $Y$ is represented by
\begin{equation}
 a=\alpha(t,z)+\Psi\!\left(\pm\sqrt{Y-v(t,z)},\,t,z\right).
 \label{eq:normalinverse}
\end{equation}
The linear coefficient is
\begin{equation}
 \partial_w\Psi(0,t,z)
 =\sqrt{\frac{2}{\partial_a^2\mathscr L(\alpha(t,z),t,z)}}.
 \label{eq:normalderivative}
\end{equation}
Consequently the inverse admits a convergent expansion in the flat
variable $z$ and the ramified target coordinate
$\sqrt{Y-v(t,z)}$, including at the completed discriminant.
\end{theorem}
\begin{proof}
Taylor's theorem with an analytic remainder gives
\begin{equation}
 \mathscr L(a,t,z)-v(t,z)
 =(a-\alpha(t,z))^2 U(a,t,z),
 \label{eq:factorfold}
\end{equation}
where
\[
 U(a,t,z)=\int_0^1(1-r)
  \partial_a^2\mathscr L\big(\alpha+r(a-\alpha),t,z\big)\dd r.
\]
At the center, $U=\partial_a^2\mathscr L/2$, bounded away from zero.
Choose an analytic square root of $U$ on a sufficiently small common
neighborhood. The coordinate
$w=(a-\alpha)\sqrt{U(a,t,z)}$ has nonzero derivative in $a$, so its
local inverse has the form $a=\alpha+\Psi(w,t,z)$, holomorphic in
$(w,z)$. Equation \eqref{eq:factorfold} becomes $Y-v=w^2$, giving
\eqref{eq:normalinverse}. Uniform derivative and boundary bounds give
common neighborhoods for small sectorial $t$. Differentiating at the
center proves \eqref{eq:normalderivative}.
\end{proof}

This is the analytic quadratic normal form with parameter. Its
classical mechanism is not the new ingredient. The information supplied
by \cref{thm:critical} is the exact $q$-product completion of its center
and discriminant. A normal form centered at $v_0$ rather than $v(t,Q)$
would be incorrect in the critical layer.

\subsection{How Stokes transport acts on the discriminant}
Let $\cL_\pm$ denote the lateral $q$-gamma carriers. Their difference is
\begin{equation}
 \cL_+(a,t)-\cL_-(a,t)
 =2\ii\sum_{n\geq1}
       \left(\sum_{d\mid n}\frac{\sin(2\pi da)}d\right)Q^n.
 \label{eq:qgammajump}
\end{equation}
The upper carrier is completed by $D_1-D_a^-$; the lower carrier is
completed by $D_1-D_a^+$. Both completed families are exactly
$\cL$. Thus the actual functions $\alpha(t,Q)$ and $v(t,Q)$ are
independent of this choice of lateral representation.

Before completion, the carrier critical values differ at leading
order by $2\ii\sin(2\pi\alpha_0)Q$. Away from the discriminant the
inverse jump is described by \cref{thm:inverse}. In the layer
$Y-v_0=O(Q)$ the derivative is $O(Q^{1/2})$, so a division by that
derivative is not a uniform approximation. The exact difference is
instead governed by two square-root expressions with shifted centers.
The elementary model
\[
 F_\pm(u)=u^2\pm\delta,\qquad
 h_\pm(Y)=\pm\sqrt{Y\mp\delta}
\]
shows why expanding in $\delta/Y$ is invalid when $Y=O(\delta)$.
The completed normal form \eqref{eq:normalinverse} is the replacement
for that expansion.

A loop around $Y=v(t,Q)$ exchanges the two inverse sheets. This is
ramification monodromy, not a new Borel Stokes jump. Similarly,
$\Re(1/t)=0$ is a loss-of-decay boundary for $Q$, not the critical
value in the target. The three phenomena must be recorded separately
in any complex inverse atlas.

\section{Applications, extensions, and genuine boundaries}
\label{sec:scope}

\subsection{Finite products and arithmetic progressions}
The theorem applies directly to
$\log(q^k;q^N)_\infty$ near $q=1$ by using $a=k/N$ and replacing $t$
with $Nt$. Finite linear combinations, ratios of infinite products,
and normalized $q$-gamma products inherit the explicit completion by
addition. Whenever their forward derivative is nonzero, the general
inverse transport theorem then supplies their completed inverse. A
zero derivative requires the actual local multiplicity and its
parameter unfolding to be examined separately.

Finite $q$-Pochhammer products and Gaussian polynomials are different.
They are algebraic or analytic functions of the parameters on their
local domains. Their inverse expansions at ordinary critical points
are convergent Puiseux expansions; the mere presence of a branch point
does not create a factorially divergent Borel tail. Stokes phenomena
enter when an infinite product, a divergent asymptotic limit, or a
nonuniform degree limit is involved. This distinguishes a finite
polynomial ramification calculation from the $q$-gamma analysis above.

\subsection{What is proved at roots of unity?}
For $q=\zeta_m\e^{-t}$, the exact factorization
\begin{equation}
 (x;q)_\infty=\prod_{r=0}^{m-1}(xq^r;q^m)_\infty,
 \qquad q^m=\e^{-mt},
 \label{eq:cyclotomicfactor}
\end{equation}
reduces the product to finitely many factors with positive-radial base.
However, their first arguments generally include nontrivial phases.
They are not all of the form $\e^{-amt}$ with a fixed real $a$.
Accordingly, \cref{thm:completion} alone is not a theorem for every
cyclotomic cusp.

There is a direct root-of-unity application when the factors depend
only on $q^m$ and have specified radial shifts:
\[
 \log\big(\e^{-am t};\e^{-mt}\big)_\infty=G_a(mt).
\]
Its flat scale is $\exp(-4\pi^2/(mt))$, and all our completion and
inverse formulas apply after that rescaling. For general phased
factors in \eqref{eq:cyclotomicfactor}, one needs the corresponding
fixed-argument Borel data, branch constants, and any moving-parameter
uniformity. The rational-shift identities in \cref{cor:fr} supply a
useful part of this structure but do not eliminate those hypotheses.

\subsection{Endpoint and coalescing-action limitations}
The constants in the basic uniform expansion are not uniform as
$a\to0$, where the product develops a zero and $\log\Gamma(a)$ is
singular. Joint limits such as $a\asymp t$ or $a\asymp Q$ require a
new local core. Likewise the normal form proved here is a stable
simple fold: the second derivative remains nonzero. An arbitrary
perturbation of an $m$-fold point for $m>2$ generally splits it into
several critical points, so replacing $1/2$ by $1/m$ without an
unfolding analysis is not justified.

All claims about exponential smallness are sectorial. We do not
assert analytic continuation of the exact infinite product through
an entire arc of $|q|=1$, nor a global inverse on the complex plane.
The Borel transform can admit continuations that the original
$q$-product does not admit on a corresponding global domain. An
inverse atlas must preserve that distinction.

\section{A general completion-sensitive ramification principle}
\label{sec:ramified}

The $q$-gamma fold suggests a useful general rule, but only after its
nondegeneracy hypothesis is made explicit. The following is a local
analytic statement underlying a class of ramified transseries.

\begin{theorem}[Action selection at a finite-multiplicity inverse]
\label{thm:ramified}
Let $F_0$ be holomorphic near $\alpha$, with
\[
 F_0(\alpha)=v_0,\quad F_0^{(j)}(\alpha)=0\ (1\leq j<m),\quad
 c_m=F_0^{(m)}(\alpha)/m!\ne0,
\]
where $m\geq2$. Let $\mathcal P(x,z)$ be holomorphic near $(\alpha,0)$
and have the form
\begin{equation}
 \mathcal P(x,z)=F_0(x)+z^r p_r(x)+z^{r+1}p_{r+1}(x)+\cdots,
 \qquad r\geq1,
 \label{eq:generalperturbation}
\end{equation}
with a convergent series. For $\eta$ in a compact set avoiding
$p_r(\alpha)$, the equation
\begin{equation}
 \mathcal P(x,z)=v_0+z^r\eta
 \label{eq:generallevel}
\end{equation}
has $m$ local solutions of the form
\begin{equation}
 x_j=\alpha+z^{r/m}\left\{u_j+O(z^{1/m})\right\},\qquad
 u_j^m=\frac{\eta-p_r(\alpha)}{c_m},
 \label{eq:generalramified}
\end{equation}
with convergent series on the cover $z=s^m$. The result is uniform for
families with uniform holomorphic bounds, $c_m$ bounded away from
zero, and the stated separation of $\eta$. Substituting
$z=\e^{-A/t}$ gives leading inverse action $rA/m$.
\end{theorem}
\begin{proof}
Set $z=s^m$ and $x=\alpha+s^r u$. Divide
\eqref{eq:generallevel} by $s^{mr}$. The quotient extends holomorphically
to $s=0$, because every higher carrier term has an additional power
$s^r$ and every later completion sector an additional power $s^m$.
At $s=0$ it is $c_mu^m+p_r(\alpha)-\eta$. Its $m$ roots are nonzero
and simple. The analytic implicit-function theorem gives the $m$
convergent branches in $s$, uniformly on compact parameter sets.
\end{proof}

This proof is an elementary Newton--Puiseux rescaling. The useful
information is the \emph{selected} exponent $r$, which must come from
the exact completion, not be guessed from the first Borel pole. At a
regular inverse ($m=1$), the same principle preserves the first
surviving action. At the nondegenerate $q$-gamma critical layer,
$m=2$ and $r=1$. Finite phase cancellation, as in
\cref{prop:delay}, can change $r$ before ramification takes place.

\begin{remark}[Why a constant-term test alone fails]
The hypotheses require the full coefficient functions of orders
below $r$ to vanish. If only their values at the critical center
vanish, lower-order derivative terms can control the inverse. For
example,
\begin{equation}
 x^2+zx+z^3=0
 \label{eq:newtoncounter}
\end{equation}
has roots
\[
 x_1=-z^2+O(z^3),\qquad x_2=-z+z^2+O(z^3),
\]
not two roots of order $z^{3/2}$. Although the completion at $x=0$
first appears as $z^3$, the term $zx$ changes the Newton polygon.
The two balances are $2\beta=1+\beta$ and $1+\beta=3$, giving
$\beta=1,2$. This counterexample is directly relevant to parameter
cancellations in $q$-analogs: vanishing of a Stokes constant or a
completion coefficient at one point is not enough to infer the
ramification scale.
\end{remark}

A full multi-action theory would replace the single variable $z$ by
an action-filtered algebra and analyze the initial equation in each
valuation. The present theorem proves the rank-one analytic case,
including uniformity and convergent ramification, without silently
assuming a Newton polygon exists for every possible transseries
support.

\section{Constructive all-orders recursions and certification}
\label{sec:algorithms}

\subsection{Regular inverse sectors}
The convergent implicit proof in \cref{thm:cusp} gives a direct
coefficient algorithm. Let
$\delta_{<n}=\sum_{j=1}^{n-1}\Delta_jz^j$. Then
\begin{equation}
 \Delta_n=-\frac1f[z^n]\left\{
 \sum_{k=2}^{n}\frac{H_a^{(k)}(h)}{k!}\delta_{<n}^{\,k}
 +\sum_{m=1}^{n}C_m(a)z^m\exp(-Am\delta_{<n})\right\}.
 \label{eq:deltarecurrence}
\end{equation}
Only finite Taylor polynomials are needed to compute the displayed
coefficient. The divisor sums $C_m$ can be generated by a sieve in
$O(N\log N)$ elementary additions through order $N$. The remaining
cost is truncated polynomial arithmetic with derivative-valued
coefficients. Formula \eqref{eq:deltarecurrence} makes resonance
combinations automatic: contributions with the same total action are
added before the next coefficient is solved.

\subsection{Critical point and value}
Write
$\alpha(t,z)=\alpha_0+\sum_{n\geq1}b_n(t)z^n$ and
$\beta_{<n}=\sum_{j<n}b_jz^j$. Then
\begin{equation}
 b_n=-\frac1\kappa[z^n]\left\{
 \sum_{k=2}^{n}\frac{\partial_a^{k+1}\cL_0(\alpha_0,t)}{k!}
                      \beta_{<n}^{\,k}
 +\sum_{m=1}^{n}z^m e_m'(\alpha_0+\beta_{<n})\right\}.
 \label{eq:criticalrecurrence}
\end{equation}
Substitute the result into \eqref{eq:auxfamily} to compute the critical
value. The denominator $\kappa$ is uniformly separated from zero;
this is where the simple-fold hypothesis enters the algorithm.

\subsection{Inverse coefficients at the completed branch point}
For the factor $U$ in \eqref{eq:factorfold}, Lagrange inversion of
$w=(a-\alpha)\sqrt{U}$ gives
\begin{equation}
 \Psi(w,t,z)=\sum_{r\geq1}\frac{w^r}{r!}
 \left.
 \partial_a^{r-1}\left[U(a,t,z)^{-r/2}\right]
 \right|_{a=\alpha(t,z)}.
 \label{eq:PsiLagrange}
\end{equation}
This is convergent on the normal-form neighborhood, not merely a
formal Puiseux identity. Together with
\eqref{eq:criticalrecurrence}, it supplies every mixed sector in the
flat parameter and ramified target coordinate. Re-expanding the
coefficient functions in $t$ gives a further, generally asymptotic,
layer of the transseries.

\subsection{What a numerical certificate must record}
A reliable inverse calculation should retain the chosen logarithmic
branch, the lateral contour, the exact completion, a target disk or
ramified coordinate, a derivative lower bound, and the norm used for
the truncation error. Near a simple fold the regular derivative lower
bound must be replaced by the nonzero quadratic coefficient and the
completed critical value.

The formulas in this article separate three errors. Truncating a
power--logarithmic carrier creates an algebraic or optimally truncated
error. Truncating the convergent $z$-expansion creates an exponential
sector error. Evaluating the coefficients numerically creates a
floating-point error. In particular, a highly accurate formal carrier
truncation is not an exact carrier unless a summation or exact
remainder representation has been supplied.

\section{Reproducible checks and numerical evidence}
\label{sec:numerics}

The accompanying \texttt{verify.py} performs exact symbolic checks and
high-precision numerical experiments. It uses only \texttt{mpmath}
and \texttt{sympy}. Numerical evidence is not used in place of the
Mellin proof, the reflection proof, or the analytic implicit-function
arguments. The output explicitly distinguishes independent evaluation
of the completion identity from tests that use that identity to define
a carrier.

\subsection{An independent completion test}
The original $G_a$ is evaluated by its absolutely convergent
logarithmic series. The Borel median is evaluated independently by
\eqref{eq:Kei} and \eqref{eq:sumK}; it is \emph{not} defined as
$G_a-P_a-D_a^0$ in this test. To accelerate the double sum, write
\begin{equation}
 b_n(a)=\sum_{d\mid n}\frac{\sin(2\pi da)}d,
 \qquad
 K_J(x)=-\frac2\pi\sum_{j=1}^{J}\frac{(2j-1)!}{x^{2j}}.
 \label{eq:Ktrunc}
\end{equation}
The approximation is
\begin{equation}
 B_{M,J}(a,t)=\sum_{j=1}^{J}c_{2j}(a)t^{2j}
 +\sum_{n=1}^{M}b_n(a)
       \left\{K(An/t)-K_J(An/t)\right\}.
 \label{eq:acceleration}
\end{equation}
For fixed $J$ and a proper $t$-sector, the large-$n$ remainder of the
kernel is $O(n^{-2J-2})$; the subtracted infinite sums give exactly the
Bernoulli terms already inserted. This makes the remaining sum
rapidly convergent. The code compares $(M,J)=(60,10)$ and $(80,12)$.
The reported errors are observed discrepancies, not interval-arithmetic
enclosures of the infinite tails.

\begin{table}[htbp]
\centering
\caption{Independent check of $G_a-P_a-D_a^0-B_{60,10}$.
The last column is the size of the explicit completion. Values are
rounded; complete output is in \texttt{verification\_results.json}.}
\label{tab:identity}
\begin{tabular}{@{}llll@{}}
\toprule
$a$ & $t$ & Observed absolute discrepancy & $|D_a^0(t)|$\\
\midrule
$1/3$ & $0.7$ & $1.262\times10^{-58}$ & $1.606\times10^{-25}$\\
$2/5$ & $1.2$ & $1.099\times10^{-53}$ & $4.171\times10^{-15}$\\
$1/3$ & $2$ & $1.354\times10^{-48}$ & $1.338\times10^{-9}$\\
$1/3$ & $0.8+0.3\ii$ & $1.012\times10^{-56}$ & $8.122\times10^{-20}$\\
\bottomrule
\end{tabular}
\end{table}

The discrepancy agrees with the change under the indicated refinement
to the displayed precision. The half-shift test separately confirms
that the Bernoulli median is zero while the dual product is nonzero.
This is a useful check that the test is not accidentally comparing two
copies of a perturbative expression.

\subsection{Critical-point and inverse-layer checks}
For this second group of tests, the exact $q$-gamma function is
computed from its convergent product, and the carrier is evaluated as
$\cL_0=\cL-\cE$ using the already independently tested completion
identity. Exact convergent derivative sums are used for Newton
iterations. Thus these are tests of the critical and inverse formulas,
not a second independent derivation of the completion theorem.

At the carrier level $Y=v_0$, the normalized inverse separation is
compared to the coefficient in \eqref{eq:realgap}. The individual
inverse approximations retain both the $Q^{1/2}$ and $Q$ terms.

\begin{table}[htbp]
\centering
\caption{Critical-layer tests at 150 decimal working digits.
The last column gives the error of each two-term inverse approximation, divided by $Q^{3/2}$.}
\label{tab:fold}
\begin{tabular}{@{}llll@{}}
\toprule
$t$ & $\alpha_0(t)$ & $(a_+-a_-)/\sqrt Q$ & $(R_-,R_+)$\\
\midrule
$1.5$ & $1.425089164991874$ & $5.99721345624$ & $(95.5583,-95.5635)$\\
$1.0$ & $1.438674038433400$ & $5.21925419341$ & $(70.9293,-70.9293)$\\
$0.75$ & $1.444978961972934$ & $4.88119865255$ & $(60.4933,-60.4933)$\\
\bottomrule
\end{tabular}
\end{table}

The critical-point shift divided by $Q$ is, respectively,
approximately $6.7733950033$, $4.1737023873$, and $3.2672115106$,
matching $-e'(\alpha_0)/\kappa$ with the predicted next-order error.
The critical-value formula is checked through order $Q^2$. These
finite experiments are consistency checks of the asymptotic orders;
they do not establish the uniform bounds, which come from the common
analytic neighborhoods in the proofs.

\subsection{Exact symbolic checks}
The symbolic tests compare several Bernoulli coefficients with the
Taylor coefficients of the holomorphic Borel kernel, verify the
$q$-gamma shift relation coefficientwise, check the two-exponential
inverse formula and the first critical-point/value formulas, and
verify the eta-quotient cancellations in \eqref{eq:eta-flat}. They do
not prove analytic continuation or convergence. No result in this
article is labeled machine-formalized.

\section{Further research questions and proposed approaches}
\label{sec:research}

The questions below are a proposed program, not a list of theorems
implicitly claimed to have been solved.

\begin{problem}[Independent audit of the scalar completion proof]
Compare \cref{thm:completion} with existing modular transformation and
Mellin-integral formulas for shifted products, then produce a precise
priority and equivalence map. The concrete deliverable is a
normalization table including the missing zero-th factors, the branch
of $\log t$, and the orientation of each lateral contour. An
independent expert check of \eqref{eq:functionaleq} and the
Mellin--principal-value identification is the highest-value next
verification step.
\end{problem}

\begin{problem}[All cyclotomic cusp factors]
Extend the completion theorem from radial shifted factors to every
phased factor in \eqref{eq:cyclotomicfactor}, uniformly in parameters
away from zeros. The target is a complete residue-class formula for
both the jump and the jump-free part, followed by an inverse theorem
with explicit branch constants. A decomposition into finitely many
fixed-argument and shifted-argument kernels is a natural starting
point; simply applying the radial theorem to all factors is not valid.
\end{problem}

\begin{problem}[Joint endpoint scaling]
Study $a\to0$ simultaneously with $t\to0$, especially $a/t$ bounded
or $a/Q$ bounded. The endpoint zero and the ordinary gamma pole should
be absorbed into a new exact core before inversion. A useful theorem
would give one uniform representation covering the fixed-$a$ formula
and the endpoint layer, with an explicit transition of Stokes and
completion coefficients.
\end{problem}

\begin{problem}[Critical cancellation and Newton polygons]
Classify parameter values of normalized $q$-products for which the
first completion coefficient vanishes at a critical point while its
parameter derivative does not. The model \eqref{eq:newtoncounter}
shows that the inverse action can then split into several valuations.
One goal is an algorithm that takes the first nonzero derivative jets
of the completion and returns the valid ramified scales and their
uniform domains.
\end{problem}

\begin{problem}[Higher-multiplicity unfoldings]
For a family with a tuned $m$-fold critical point, determine the
completed discriminant and all nearby simple critical points after
exponentially small perturbation. The level theorem
\cref{thm:ramified} is not a classification of that discriminant.
Weierstrass preparation in the inverse variable, combined with the
exact $Q$-series, should reduce the problem to a polynomial unfolding
whose coefficients are sectorial transseries.
\end{problem}

\begin{problem}[Borel singularities of the critical coefficients]
Determine the analytic continuation and alien derivatives of the
formal $t$-expansions of $\alpha_0$, $v_0$, and the coefficients of
$\Psi$. Resurgent implicit-function closure is a starting point, but
it does not by itself provide a sharp list of singularities or their
local types. A concrete first target is the leading singularity of
the critical-point series, computed from the $a$-derivative of the
explicit Borel kernel and the implicit critical equation.
\end{problem}

\begin{problem}[Sharp convergence radii in the exponential variable]
The proofs give common positive radii in $z$, but not the exact
nearest singularity of $\alpha(t,z)$ or the regular inverse series.
Determine whether the first obstruction is a collision of critical
points, a zero of a dual product factor, or escape from the chosen
inverse chart. This is an analytic continuation problem, not a
question answered by a finite list of coefficients.
\end{problem}

\begin{problem}[Infinite reflection weights]
Replace the finite phase measure by a countable or continuous measure
on the circle. Under what summability conditions do the two reflection
kernels remain exactly the even and odd subspaces? Fourier uniqueness
suggests a positive answer for finite measures, but interchange with
the Borel and dual-product transforms must be justified. The
finite-phase action-delay bound has no immediate infinite analogue.
\end{problem}

\begin{problem}[Several actions and resonance blocks]
Develop a uniform version of \cref{thm:ramified} for several complex
actions and for families in which their real parts exchange order.
The correct input is an action filtration with locally finite sums,
not an arbitrary list of exponentially small terms. The target is a
branch-compatible inverse cocycle whose sector regrouping remains
valid when actions resonate or coalesce.
\end{problem}

\begin{problem}[Character-twisted Lambert completions]
Determine which parts of the Mellin proof extend to the broader
character-twisted Lambert families of \cite{bd}. The decisive issue is
whether the functional equation splits the Mellin kernel into a
Laplace-summed perturbative part and an explicit convergent flat part
with the proposed phases. The present scalar shifted-product theorem
does not establish every such phase prescription.
\end{problem}

\begin{problem}[Nonlinear vector-valued modular transport]
For vectors of weighted $q$-products, identify conditions under which
an analytically defined inverse observable inherits a compatible
modular or quantum-modular transition law. Linear antisymmetrization
can eliminate a completion, but nonlinear inversion creates mixed
sectors. An exact transition map such as \eqref{eq:cocycle}, rather
than coordinatewise arithmetic averaging, should be the basic object.
\end{problem}

\begin{problem}[Formalization and verified numerics]
Formalize first the finite-weight Vandermonde criterion, the filtered
fixed-point lemma, and the algebraic coefficient recursions. In
parallel, implement interval-enclosed evaluation of the scalar
principal-value kernel and derivative bounds on the fold polydisk.
A full formal proof of the Mellin contour deformation will require a
substantial complex-analysis library; numerical agreement alone
should never be presented as that formal proof.
\end{problem}

\section{Conclusions}
The exact completion, rather than the formal tail alone, is the
natural input for complex transseries reversion. In the shifted
$q$-product family, a Mellin functional equation identifies both the
Stokes-sensitive sine part and the jump-free cosine part. This proves
a scalar completion identity appearing as a recent conjecture and
leads to a finite-weight reflection classification.

Regular inverse transport then converts the completed $q=1$ cusp into
a power--logarithmic expansion dressed by convergent exponential
sectors. At the $q$-gamma minimum, the critical value itself moves by
order $Q$, and the inverse branches respond at order $Q^{1/2}$. The
completed quadratic normal form remains valid where regular inverse
estimates fail. Its location, not merely its existence, is determined
by the exact completion theorem.

The resulting theory is a concrete part of a broader complex
transseries inverse atlas: formal support, sectorial summation,
completion selection, and sheet monodromy remain separate data that
must be made compatible. The article supplies that compatibility for
the stated shifted-product and $q$-gamma charts, rather than claiming
a single global inversion theory for every possible transseries.

\appendix
\section{Residue audit and finite-order coefficients}
\label{app:residues}
For reference, the residues in \eqref{eq:mellinG} can be checked
without any Borel argument. At $s=-2n$, $n\geq1$,
\[
 \Res\Gamma(s)=\frac1{(2n)!},\qquad
 \zeta(1-2n)=-\frac{B_{2n}}{2n},\qquad
 \zeta(-2n,a)=-\frac{B_{2n+1}(a)}{2n+1}.
\]
The overall minus sign in \eqref{eq:mellinG} therefore gives exactly
$c_{2n}(a)$ in \eqref{eq:tail}. At the negative odd integers below
$-1$, the zero of $\zeta(s+1)$ cancels the gamma pole. At $s=-1$ the
three factors give the linear coefficient $B_2(a)/4$.

The Stokes residue sign is independent of this Mellin residue sign.
For a contour going from left to right above a positive pole, the
small upper arc is clockwise, and contributes $-\pi\ii$ times the
residue. The lower arc contributes $+\pi\ii$ times the residue.
Their difference is $-2\pi\ii$ times the residue, as used in
\eqref{eq:sign}. The exact $G_a$ pairs its upper Borel sum with the
\emph{negative-phase} dual product, as in
\eqref{eq:completionplus}. Reversing that pairing would change the
function by a nonzero exponential jump.

The exact $q$-gamma shift identity is
\begin{equation}
 \cL(a+1,t)-\cL(a,t)
 =\log(1-\e^{-at})-\log(1-\e^{-t}).
 \label{eq:qgammashift}
\end{equation}
The linear coefficients in \eqref{eq:ell1} have difference $(1-a)/2$,
and the even coefficients obey
\begin{equation}
 \ell_{2n}(a+1)-\ell_{2n}(a)
 =\frac{B_{2n}(a^{2n}-1)}{2n(2n)!}.
 \label{eq:ellshift}
\end{equation}
The completion is periodic in $a$, so it cancels in
\eqref{eq:qgammashift}. These identities provide independent
coefficient checks with exact rational arithmetic.

\section{Internal proof review and statement boundaries}
\label{app:audit}
This is an internal review checklist, not an independent referee
report. The main points audited are the following.

\paragraph{Mellin strip.}
The shifted line is $1<c<2$ in the $u$-variable. The lower bound makes
$\zeta(u)$ and the double Dirichlet sums absolutely convergent; the
upper bound places the scalar principal-value kernel in its Mellin
strip. Choosing an arbitrary positive line would miss one of these
requirements.

\paragraph{Principal values.}
The infinite Borel-kernel integral is not handled by an unqualified
Tonelli argument. The proof integrates on a nonreal ray, obtains an
absolute $k^{-3}m^{-2}$ majorant, deforms each term, and uses the
Mellin bound to sum the resulting principal-value kernels. The scalar
Mellin transform itself uses subtraction at its one pole.

\paragraph{Shift analyticity.}
The real Fourier representation is not used for nonreal $a$. The
sine-quotient kernel supplies holomorphic continuation, and the
$a\mapsto a+1$ relation supplies the convergent tail correction needed
near the gamma minimum.

\paragraph{Finite versus infinite weights.}
The two reflection equivalences and the delay bound use a finite
Vandermonde system. No result for unbounded or arbitrary distributional
weights is inferred from that proof.

\paragraph{Completion and median.}
The term $D_a^0$ is not equated with half a Stokes jump. Its cosine
coefficients are independent of the sine coefficients that govern
the jump. The arithmetic median is used only in the additive
logarithmic equation. Nonlinear inverse reconstruction is performed
with an analytic inverse theorem.

\paragraph{Uniformity and ramification.}
The regular inverse formulas use disks with derivative control. The
critical formulas promote $Q$ to an independent $z$ and use common
polydisks. The layer formula excludes its scaled branch point; the
completed normal form covers that point. The general action-selection
principle excludes a degenerate initial polynomial, and
\eqref{eq:newtoncounter} records a case where a naive scale rule fails.

\paragraph{What has not been claimed.}
The article does not prove a global inverse for all complex
transseries, a full cyclotomic-cusp classification, an arbitrary
higher-critical-point unfolding theorem, or the general conjectural
phase formulas for character-twisted Lambert series. It does prove
the explicit scalar completion identity and the stated inverse
consequences. Its numerical checks are neither interval proofs nor
Lean proofs, and its source audit does not certify first-publication
priority.

\section{Source snapshot and reproducibility}
\label{app:source}
The repository snapshot used for comparison is the commit printed in
\cref{sec:introduction}; the fetched nonlinear-transport article was
read at that commit, rather than inferred solely from a search
snippet. The principal external status comparison uses
Fantini--Rella, arXiv:2506.08265v2, dated 1 April 2026, where the scalar
identity is explicitly labeled Conjecture~5. The later works listed
in the bibliography were considered to avoid relying only on a 2025
literature snapshot. The comparison does not establish that no
equivalent identity exists in older transformation-formula literature.

The distributed package contains the standalone \LaTeX{} source,
its compiled PDF, \texttt{verify.py}, a captured numerical-results
JSON file, \texttt{requirements.txt}, \texttt{README.md}, and a compact
proof-review note. Compilation requires a standard \TeX{} installation
with the packages listed in the source. Two or more \texttt{pdflatex}
passes resolve cross-references. The bibliography is embedded, so
Bib\TeX{} is not required.

Run the numerical verification with
\begin{quote}\small\ttfamily
python verify.py --dps 150 --output verification\_results.json
\end{quote}
The tests use decimal strings or exact rationals for their inputs.
No binary machine-precision approximation of $\pi$ is reused in a
high-precision calculation. Changing the precision can alter the
last displayed digits and the noise floor in identically vanishing
tests without changing the mathematical assertions.

\begin{thebibliography}{99}

\bibitem{repo}
V. Reshetnikov, \emph{ProveIt}, repository snapshot
\texttt{b484725c893a7a9bde856151988996aef6965cb0}, inspected
4 October 2026.\newline
\url{https://github.com/VladimirReshetnikov/ProveIt/tree/b484725c893a7a9bde856151988996aef6965cb0}.
The consolidated transseries group is under
\path{Analysis/Transseries/docs/series-and-transseries/}.

\bibitem{repotransport}
ProveIt research report, \emph{Nonlinear Stokes Transport under
Logarithmic Inversion}, 29 September 2026, source inspected at the
snapshot in \cite{repo}.\newline
\path{Analysis/Transseries/docs/series-and-transseries/Nonlinear_Stokes_Transport_Logarithmic_Inversion/article.tex}.

\bibitem{repocrossover}
ProveIt research report, \emph{Uniform Resurgent Crossover for Gaussian
Binomials}, source group inspected in \cite{repo}.\newline
\path{Analysis/Transseries/docs/series-and-transseries/Uniform_Resurgent_Crossover_Gaussian_Binomials/}.

\bibitem{fr}
V. Fantini and C. Rella,
\emph{Modular resurgence, $q$-Pochhammer symbols, and quantum operators
from mirror curves}, arXiv:2506.08265v2, 1 April 2026.
In particular, Section 3.1.3, Conjecture 5, and equations (60)--(62).\newline
\url{https://arxiv.org/abs/2506.08265v2}.

\bibitem{vectors}
M. C. N. Cheng, I. Coman, V. Fantini, and C. Rella,
\emph{Modular resurgent structures for vectors},
arXiv:2608.08902v1, 9 August 2026.\newline
\url{https://arxiv.org/abs/2608.08902v1}.

\bibitem{bd}
D. Broadhurst and D. Dorigoni,
\emph{Resurgent Lambert series with characters},
arXiv:2507.21352v2, 30 July 2026.\newline
\url{https://arxiv.org/abs/2507.21352v2}.

\bibitem{sauzin}
D. Sauzin, \emph{Nonlinear analysis with resurgent functions},
arXiv:1212.4477.\newline
\url{https://arxiv.org/abs/1212.4477}.

\bibitem{sauzinintro}
D. Sauzin, \emph{Introduction to 1-summability and resurgence},
arXiv:1405.0356.\newline
\url{https://arxiv.org/abs/1405.0356}.

\bibitem{gk}
S. Garoufalidis and R. Kashaev,
\emph{Resurgence of Faddeev's quantum dilogarithm},
IRMA Lectures in Mathematics and Theoretical Physics \textbf{33}
(2021), 257--272; arXiv:2008.12465.\newline
\url{https://arxiv.org/abs/2008.12465}.

\bibitem{dlmfzeta}
NIST Digital Library of Mathematical Functions,
\emph{Section 25.11: Hurwitz Zeta Function}, including the functional
equation, Bernoulli-polynomial values, and the derivative at zero;
consulted 4 October 2026.\newline
\url{https://dlmf.nist.gov/25.11}.

\bibitem{dlmfqgamma}
NIST Digital Library of Mathematical Functions,
\emph{Section 5.18: $q$-Gamma and $q$-Beta Functions},
consulted 4 October 2026.\newline
\url{https://dlmf.nist.gov/5.18}.

\end{thebibliography}
\end{document}
